\documentclass[11pt]{article}

    \usepackage[T2A]{fontenc}
\usepackage[utf8]{inputenc}
\usepackage[english, russian]{babel}

\usepackage[breakable]{tcolorbox}
    \usepackage{parskip} % Stop auto-indenting (to mimic markdown behaviour)
    
    \usepackage{iftex}
    \ifPDFTeX
    	\usepackage[T1]{fontenc}
    	\usepackage{mathpazo}
    \else
    	\usepackage{fontspec}
    \fi

    % Basic figure setup, for now with no caption control since it's done
    % automatically by Pandoc (which extracts ![](path) syntax from Markdown).
    \usepackage{graphicx}
    % Maintain compatibility with old templates. Remove in nbconvert 6.0
    \let\Oldincludegraphics\includegraphics
    % Ensure that by default, figures have no caption (until we provide a
    % proper Figure object with a Caption API and a way to capture that
    % in the conversion process - todo).
    \usepackage{caption}
    \DeclareCaptionFormat{nocaption}{}
    \captionsetup{format=nocaption,aboveskip=0pt,belowskip=0pt}

    \usepackage{float}
    \floatplacement{figure}{H} % forces figures to be placed at the correct location
    \usepackage{xcolor} % Allow colors to be defined
    \usepackage{enumerate} % Needed for markdown enumerations to work
    \usepackage{geometry} % Used to adjust the document margins
    \usepackage{amsmath} % Equations
    \usepackage{amssymb} % Equations
    \usepackage{textcomp} % defines textquotesingle
    % Hack from http://tex.stackexchange.com/a/47451/13684:
    \AtBeginDocument{%
        \def\PYZsq{\textquotesingle}% Upright quotes in Pygmentized code
    }
    \usepackage{upquote} % Upright quotes for verbatim code
    \usepackage{eurosym} % defines \euro
    \usepackage[mathletters]{ucs} % Extended unicode (utf-8) support
    \usepackage{fancyvrb} % verbatim replacement that allows latex
    \usepackage{grffile} % extends the file name processing of package graphics 
                         % to support a larger range
    \makeatletter % fix for old versions of grffile with XeLaTeX
    \@ifpackagelater{grffile}{2019/11/01}
    {
      % Do nothing on new versions
    }
    {
      \def\Gread@@xetex#1{%
        \IfFileExists{"\Gin@base".bb}%
        {\Gread@eps{\Gin@base.bb}}%
        {\Gread@@xetex@aux#1}%
      }
    }
    \makeatother
    \usepackage[Export]{adjustbox} % Used to constrain images to a maximum size
    \adjustboxset{max size={0.9\linewidth}{0.9\paperheight}}

    % The hyperref package gives us a pdf with properly built
    % internal navigation ('pdf bookmarks' for the table of contents,
    % internal cross-reference links, web links for URLs, etc.)
    \usepackage{hyperref}
    % The default LaTeX title has an obnoxious amount of whitespace. By default,
    % titling removes some of it. It also provides customization options.
    \usepackage{titling}
    \usepackage{longtable} % longtable support required by pandoc >1.10
    \usepackage{booktabs}  % table support for pandoc > 1.12.2
    \usepackage[inline]{enumitem} % IRkernel/repr support (it uses the enumerate* environment)
    \usepackage[normalem]{ulem} % ulem is needed to support strikethroughs (\sout)
                                % normalem makes italics be italics, not underlines
    \usepackage{mathrsfs}
    
\usepackage{wrapfig}
\usepackage[rightcaption]{sidecap}
\providecommand{\keywords}[1]{\textbf{\textit{Keywords:}} #1}

\author{A. Yu. Drozdov}

    
    % Colors for the hyperref package
    \definecolor{urlcolor}{rgb}{0,.145,.698}
    \definecolor{linkcolor}{rgb}{.71,0.21,0.01}
    \definecolor{citecolor}{rgb}{.12,.54,.11}

    % ANSI colors
    \definecolor{ansi-black}{HTML}{3E424D}
    \definecolor{ansi-black-intense}{HTML}{282C36}
    \definecolor{ansi-red}{HTML}{E75C58}
    \definecolor{ansi-red-intense}{HTML}{B22B31}
    \definecolor{ansi-green}{HTML}{00A250}
    \definecolor{ansi-green-intense}{HTML}{007427}
    \definecolor{ansi-yellow}{HTML}{DDB62B}
    \definecolor{ansi-yellow-intense}{HTML}{B27D12}
    \definecolor{ansi-blue}{HTML}{208FFB}
    \definecolor{ansi-blue-intense}{HTML}{0065CA}
    \definecolor{ansi-magenta}{HTML}{D160C4}
    \definecolor{ansi-magenta-intense}{HTML}{A03196}
    \definecolor{ansi-cyan}{HTML}{60C6C8}
    \definecolor{ansi-cyan-intense}{HTML}{258F8F}
    \definecolor{ansi-white}{HTML}{C5C1B4}
    \definecolor{ansi-white-intense}{HTML}{A1A6B2}
    \definecolor{ansi-default-inverse-fg}{HTML}{FFFFFF}
    \definecolor{ansi-default-inverse-bg}{HTML}{000000}

    % common color for the border for error outputs.
    \definecolor{outerrorbackground}{HTML}{FFDFDF}

    % commands and environments needed by pandoc snippets
    % extracted from the output of `pandoc -s`
    \providecommand{\tightlist}{%
      \setlength{\itemsep}{0pt}\setlength{\parskip}{0pt}}
    \DefineVerbatimEnvironment{Highlighting}{Verbatim}{commandchars=\\\{\}}
    % Add ',fontsize=\small' for more characters per line
    \newenvironment{Shaded}{}{}
    \newcommand{\KeywordTok}[1]{\textcolor[rgb]{0.00,0.44,0.13}{\textbf{{#1}}}}
    \newcommand{\DataTypeTok}[1]{\textcolor[rgb]{0.56,0.13,0.00}{{#1}}}
    \newcommand{\DecValTok}[1]{\textcolor[rgb]{0.25,0.63,0.44}{{#1}}}
    \newcommand{\BaseNTok}[1]{\textcolor[rgb]{0.25,0.63,0.44}{{#1}}}
    \newcommand{\FloatTok}[1]{\textcolor[rgb]{0.25,0.63,0.44}{{#1}}}
    \newcommand{\CharTok}[1]{\textcolor[rgb]{0.25,0.44,0.63}{{#1}}}
    \newcommand{\StringTok}[1]{\textcolor[rgb]{0.25,0.44,0.63}{{#1}}}
    \newcommand{\CommentTok}[1]{\textcolor[rgb]{0.38,0.63,0.69}{\textit{{#1}}}}
    \newcommand{\OtherTok}[1]{\textcolor[rgb]{0.00,0.44,0.13}{{#1}}}
    \newcommand{\AlertTok}[1]{\textcolor[rgb]{1.00,0.00,0.00}{\textbf{{#1}}}}
    \newcommand{\FunctionTok}[1]{\textcolor[rgb]{0.02,0.16,0.49}{{#1}}}
    \newcommand{\RegionMarkerTok}[1]{{#1}}
    \newcommand{\ErrorTok}[1]{\textcolor[rgb]{1.00,0.00,0.00}{\textbf{{#1}}}}
    \newcommand{\NormalTok}[1]{{#1}}
    
    % Additional commands for more recent versions of Pandoc
    \newcommand{\ConstantTok}[1]{\textcolor[rgb]{0.53,0.00,0.00}{{#1}}}
    \newcommand{\SpecialCharTok}[1]{\textcolor[rgb]{0.25,0.44,0.63}{{#1}}}
    \newcommand{\VerbatimStringTok}[1]{\textcolor[rgb]{0.25,0.44,0.63}{{#1}}}
    \newcommand{\SpecialStringTok}[1]{\textcolor[rgb]{0.73,0.40,0.53}{{#1}}}
    \newcommand{\ImportTok}[1]{{#1}}
    \newcommand{\DocumentationTok}[1]{\textcolor[rgb]{0.73,0.13,0.13}{\textit{{#1}}}}
    \newcommand{\AnnotationTok}[1]{\textcolor[rgb]{0.38,0.63,0.69}{\textbf{\textit{{#1}}}}}
    \newcommand{\CommentVarTok}[1]{\textcolor[rgb]{0.38,0.63,0.69}{\textbf{\textit{{#1}}}}}
    \newcommand{\VariableTok}[1]{\textcolor[rgb]{0.10,0.09,0.49}{{#1}}}
    \newcommand{\ControlFlowTok}[1]{\textcolor[rgb]{0.00,0.44,0.13}{\textbf{{#1}}}}
    \newcommand{\OperatorTok}[1]{\textcolor[rgb]{0.40,0.40,0.40}{{#1}}}
    \newcommand{\BuiltInTok}[1]{{#1}}
    \newcommand{\ExtensionTok}[1]{{#1}}
    \newcommand{\PreprocessorTok}[1]{\textcolor[rgb]{0.74,0.48,0.00}{{#1}}}
    \newcommand{\AttributeTok}[1]{\textcolor[rgb]{0.49,0.56,0.16}{{#1}}}
    \newcommand{\InformationTok}[1]{\textcolor[rgb]{0.38,0.63,0.69}{\textbf{\textit{{#1}}}}}
    \newcommand{\WarningTok}[1]{\textcolor[rgb]{0.38,0.63,0.69}{\textbf{\textit{{#1}}}}}
    
    
    % Define a nice break command that doesn't care if a line doesn't already
    % exist.
    \def\br{\hspace*{\fill} \\* }
    % Math Jax compatibility definitions
    \def\gt{>}
    \def\lt{<}
    \let\Oldtex\TeX
    \let\Oldlatex\LaTeX
    \renewcommand{\TeX}{\textrm{\Oldtex}}
    \renewcommand{\LaTeX}{\textrm{\Oldlatex}}
    % Document parameters
    % Document title
    \title{MenDrive\_ponderomotive}
    
    
    
    
    
% Pygments definitions
\makeatletter
\def\PY@reset{\let\PY@it=\relax \let\PY@bf=\relax%
    \let\PY@ul=\relax \let\PY@tc=\relax%
    \let\PY@bc=\relax \let\PY@ff=\relax}
\def\PY@tok#1{\csname PY@tok@#1\endcsname}
\def\PY@toks#1+{\ifx\relax#1\empty\else%
    \PY@tok{#1}\expandafter\PY@toks\fi}
\def\PY@do#1{\PY@bc{\PY@tc{\PY@ul{%
    \PY@it{\PY@bf{\PY@ff{#1}}}}}}}
\def\PY#1#2{\PY@reset\PY@toks#1+\relax+\PY@do{#2}}

\@namedef{PY@tok@w}{\def\PY@tc##1{\textcolor[rgb]{0.73,0.73,0.73}{##1}}}
\@namedef{PY@tok@c}{\let\PY@it=\textit\def\PY@tc##1{\textcolor[rgb]{0.25,0.50,0.50}{##1}}}
\@namedef{PY@tok@cp}{\def\PY@tc##1{\textcolor[rgb]{0.74,0.48,0.00}{##1}}}
\@namedef{PY@tok@k}{\let\PY@bf=\textbf\def\PY@tc##1{\textcolor[rgb]{0.00,0.50,0.00}{##1}}}
\@namedef{PY@tok@kp}{\def\PY@tc##1{\textcolor[rgb]{0.00,0.50,0.00}{##1}}}
\@namedef{PY@tok@kt}{\def\PY@tc##1{\textcolor[rgb]{0.69,0.00,0.25}{##1}}}
\@namedef{PY@tok@o}{\def\PY@tc##1{\textcolor[rgb]{0.40,0.40,0.40}{##1}}}
\@namedef{PY@tok@ow}{\let\PY@bf=\textbf\def\PY@tc##1{\textcolor[rgb]{0.67,0.13,1.00}{##1}}}
\@namedef{PY@tok@nb}{\def\PY@tc##1{\textcolor[rgb]{0.00,0.50,0.00}{##1}}}
\@namedef{PY@tok@nf}{\def\PY@tc##1{\textcolor[rgb]{0.00,0.00,1.00}{##1}}}
\@namedef{PY@tok@nc}{\let\PY@bf=\textbf\def\PY@tc##1{\textcolor[rgb]{0.00,0.00,1.00}{##1}}}
\@namedef{PY@tok@nn}{\let\PY@bf=\textbf\def\PY@tc##1{\textcolor[rgb]{0.00,0.00,1.00}{##1}}}
\@namedef{PY@tok@ne}{\let\PY@bf=\textbf\def\PY@tc##1{\textcolor[rgb]{0.82,0.25,0.23}{##1}}}
\@namedef{PY@tok@nv}{\def\PY@tc##1{\textcolor[rgb]{0.10,0.09,0.49}{##1}}}
\@namedef{PY@tok@no}{\def\PY@tc##1{\textcolor[rgb]{0.53,0.00,0.00}{##1}}}
\@namedef{PY@tok@nl}{\def\PY@tc##1{\textcolor[rgb]{0.63,0.63,0.00}{##1}}}
\@namedef{PY@tok@ni}{\let\PY@bf=\textbf\def\PY@tc##1{\textcolor[rgb]{0.60,0.60,0.60}{##1}}}
\@namedef{PY@tok@na}{\def\PY@tc##1{\textcolor[rgb]{0.49,0.56,0.16}{##1}}}
\@namedef{PY@tok@nt}{\let\PY@bf=\textbf\def\PY@tc##1{\textcolor[rgb]{0.00,0.50,0.00}{##1}}}
\@namedef{PY@tok@nd}{\def\PY@tc##1{\textcolor[rgb]{0.67,0.13,1.00}{##1}}}
\@namedef{PY@tok@s}{\def\PY@tc##1{\textcolor[rgb]{0.73,0.13,0.13}{##1}}}
\@namedef{PY@tok@sd}{\let\PY@it=\textit\def\PY@tc##1{\textcolor[rgb]{0.73,0.13,0.13}{##1}}}
\@namedef{PY@tok@si}{\let\PY@bf=\textbf\def\PY@tc##1{\textcolor[rgb]{0.73,0.40,0.53}{##1}}}
\@namedef{PY@tok@se}{\let\PY@bf=\textbf\def\PY@tc##1{\textcolor[rgb]{0.73,0.40,0.13}{##1}}}
\@namedef{PY@tok@sr}{\def\PY@tc##1{\textcolor[rgb]{0.73,0.40,0.53}{##1}}}
\@namedef{PY@tok@ss}{\def\PY@tc##1{\textcolor[rgb]{0.10,0.09,0.49}{##1}}}
\@namedef{PY@tok@sx}{\def\PY@tc##1{\textcolor[rgb]{0.00,0.50,0.00}{##1}}}
\@namedef{PY@tok@m}{\def\PY@tc##1{\textcolor[rgb]{0.40,0.40,0.40}{##1}}}
\@namedef{PY@tok@gh}{\let\PY@bf=\textbf\def\PY@tc##1{\textcolor[rgb]{0.00,0.00,0.50}{##1}}}
\@namedef{PY@tok@gu}{\let\PY@bf=\textbf\def\PY@tc##1{\textcolor[rgb]{0.50,0.00,0.50}{##1}}}
\@namedef{PY@tok@gd}{\def\PY@tc##1{\textcolor[rgb]{0.63,0.00,0.00}{##1}}}
\@namedef{PY@tok@gi}{\def\PY@tc##1{\textcolor[rgb]{0.00,0.63,0.00}{##1}}}
\@namedef{PY@tok@gr}{\def\PY@tc##1{\textcolor[rgb]{1.00,0.00,0.00}{##1}}}
\@namedef{PY@tok@ge}{\let\PY@it=\textit}
\@namedef{PY@tok@gs}{\let\PY@bf=\textbf}
\@namedef{PY@tok@gp}{\let\PY@bf=\textbf\def\PY@tc##1{\textcolor[rgb]{0.00,0.00,0.50}{##1}}}
\@namedef{PY@tok@go}{\def\PY@tc##1{\textcolor[rgb]{0.53,0.53,0.53}{##1}}}
\@namedef{PY@tok@gt}{\def\PY@tc##1{\textcolor[rgb]{0.00,0.27,0.87}{##1}}}
\@namedef{PY@tok@err}{\def\PY@bc##1{{\setlength{\fboxsep}{\string -\fboxrule}\fcolorbox[rgb]{1.00,0.00,0.00}{1,1,1}{\strut ##1}}}}
\@namedef{PY@tok@kc}{\let\PY@bf=\textbf\def\PY@tc##1{\textcolor[rgb]{0.00,0.50,0.00}{##1}}}
\@namedef{PY@tok@kd}{\let\PY@bf=\textbf\def\PY@tc##1{\textcolor[rgb]{0.00,0.50,0.00}{##1}}}
\@namedef{PY@tok@kn}{\let\PY@bf=\textbf\def\PY@tc##1{\textcolor[rgb]{0.00,0.50,0.00}{##1}}}
\@namedef{PY@tok@kr}{\let\PY@bf=\textbf\def\PY@tc##1{\textcolor[rgb]{0.00,0.50,0.00}{##1}}}
\@namedef{PY@tok@bp}{\def\PY@tc##1{\textcolor[rgb]{0.00,0.50,0.00}{##1}}}
\@namedef{PY@tok@fm}{\def\PY@tc##1{\textcolor[rgb]{0.00,0.00,1.00}{##1}}}
\@namedef{PY@tok@vc}{\def\PY@tc##1{\textcolor[rgb]{0.10,0.09,0.49}{##1}}}
\@namedef{PY@tok@vg}{\def\PY@tc##1{\textcolor[rgb]{0.10,0.09,0.49}{##1}}}
\@namedef{PY@tok@vi}{\def\PY@tc##1{\textcolor[rgb]{0.10,0.09,0.49}{##1}}}
\@namedef{PY@tok@vm}{\def\PY@tc##1{\textcolor[rgb]{0.10,0.09,0.49}{##1}}}
\@namedef{PY@tok@sa}{\def\PY@tc##1{\textcolor[rgb]{0.73,0.13,0.13}{##1}}}
\@namedef{PY@tok@sb}{\def\PY@tc##1{\textcolor[rgb]{0.73,0.13,0.13}{##1}}}
\@namedef{PY@tok@sc}{\def\PY@tc##1{\textcolor[rgb]{0.73,0.13,0.13}{##1}}}
\@namedef{PY@tok@dl}{\def\PY@tc##1{\textcolor[rgb]{0.73,0.13,0.13}{##1}}}
\@namedef{PY@tok@s2}{\def\PY@tc##1{\textcolor[rgb]{0.73,0.13,0.13}{##1}}}
\@namedef{PY@tok@sh}{\def\PY@tc##1{\textcolor[rgb]{0.73,0.13,0.13}{##1}}}
\@namedef{PY@tok@s1}{\def\PY@tc##1{\textcolor[rgb]{0.73,0.13,0.13}{##1}}}
\@namedef{PY@tok@mb}{\def\PY@tc##1{\textcolor[rgb]{0.40,0.40,0.40}{##1}}}
\@namedef{PY@tok@mf}{\def\PY@tc##1{\textcolor[rgb]{0.40,0.40,0.40}{##1}}}
\@namedef{PY@tok@mh}{\def\PY@tc##1{\textcolor[rgb]{0.40,0.40,0.40}{##1}}}
\@namedef{PY@tok@mi}{\def\PY@tc##1{\textcolor[rgb]{0.40,0.40,0.40}{##1}}}
\@namedef{PY@tok@il}{\def\PY@tc##1{\textcolor[rgb]{0.40,0.40,0.40}{##1}}}
\@namedef{PY@tok@mo}{\def\PY@tc##1{\textcolor[rgb]{0.40,0.40,0.40}{##1}}}
\@namedef{PY@tok@ch}{\let\PY@it=\textit\def\PY@tc##1{\textcolor[rgb]{0.25,0.50,0.50}{##1}}}
\@namedef{PY@tok@cm}{\let\PY@it=\textit\def\PY@tc##1{\textcolor[rgb]{0.25,0.50,0.50}{##1}}}
\@namedef{PY@tok@cpf}{\let\PY@it=\textit\def\PY@tc##1{\textcolor[rgb]{0.25,0.50,0.50}{##1}}}
\@namedef{PY@tok@c1}{\let\PY@it=\textit\def\PY@tc##1{\textcolor[rgb]{0.25,0.50,0.50}{##1}}}
\@namedef{PY@tok@cs}{\let\PY@it=\textit\def\PY@tc##1{\textcolor[rgb]{0.25,0.50,0.50}{##1}}}

\def\PYZbs{\char`\\}
\def\PYZus{\char`\_}
\def\PYZob{\char`\{}
\def\PYZcb{\char`\}}
\def\PYZca{\char`\^}
\def\PYZam{\char`\&}
\def\PYZlt{\char`\<}
\def\PYZgt{\char`\>}
\def\PYZsh{\char`\#}
\def\PYZpc{\char`\%}
\def\PYZdl{\char`\$}
\def\PYZhy{\char`\-}
\def\PYZsq{\char`\'}
\def\PYZdq{\char`\"}
\def\PYZti{\char`\~}
% for compatibility with earlier versions
\def\PYZat{@}
\def\PYZlb{[}
\def\PYZrb{]}
\makeatother


    % For linebreaks inside Verbatim environment from package fancyvrb. 
    \makeatletter
        \newbox\Wrappedcontinuationbox 
        \newbox\Wrappedvisiblespacebox 
        \newcommand*\Wrappedvisiblespace {\textcolor{red}{\textvisiblespace}} 
        \newcommand*\Wrappedcontinuationsymbol {\textcolor{red}{\llap{\tiny$\m@th\hookrightarrow$}}} 
        \newcommand*\Wrappedcontinuationindent {3ex } 
        \newcommand*\Wrappedafterbreak {\kern\Wrappedcontinuationindent\copy\Wrappedcontinuationbox} 
        % Take advantage of the already applied Pygments mark-up to insert 
        % potential linebreaks for TeX processing. 
        %        {, <, #, %, $, ' and ": go to next line. 
        %        _, }, ^, &, >, - and ~: stay at end of broken line. 
        % Use of \textquotesingle for straight quote. 
        \newcommand*\Wrappedbreaksatspecials {% 
            \def\PYGZus{\discretionary{\char`\_}{\Wrappedafterbreak}{\char`\_}}% 
            \def\PYGZob{\discretionary{}{\Wrappedafterbreak\char`\{}{\char`\{}}% 
            \def\PYGZcb{\discretionary{\char`\}}{\Wrappedafterbreak}{\char`\}}}% 
            \def\PYGZca{\discretionary{\char`\^}{\Wrappedafterbreak}{\char`\^}}% 
            \def\PYGZam{\discretionary{\char`\&}{\Wrappedafterbreak}{\char`\&}}% 
            \def\PYGZlt{\discretionary{}{\Wrappedafterbreak\char`\<}{\char`\<}}% 
            \def\PYGZgt{\discretionary{\char`\>}{\Wrappedafterbreak}{\char`\>}}% 
            \def\PYGZsh{\discretionary{}{\Wrappedafterbreak\char`\#}{\char`\#}}% 
            \def\PYGZpc{\discretionary{}{\Wrappedafterbreak\char`\%}{\char`\%}}% 
            \def\PYGZdl{\discretionary{}{\Wrappedafterbreak\char`\$}{\char`\$}}% 
            \def\PYGZhy{\discretionary{\char`\-}{\Wrappedafterbreak}{\char`\-}}% 
            \def\PYGZsq{\discretionary{}{\Wrappedafterbreak\textquotesingle}{\textquotesingle}}% 
            \def\PYGZdq{\discretionary{}{\Wrappedafterbreak\char`\"}{\char`\"}}% 
            \def\PYGZti{\discretionary{\char`\~}{\Wrappedafterbreak}{\char`\~}}% 
        } 
        % Some characters . , ; ? ! / are not pygmentized. 
        % This macro makes them "active" and they will insert potential linebreaks 
        \newcommand*\Wrappedbreaksatpunct {% 
            \lccode`\~`\.\lowercase{\def~}{\discretionary{\hbox{\char`\.}}{\Wrappedafterbreak}{\hbox{\char`\.}}}% 
            \lccode`\~`\,\lowercase{\def~}{\discretionary{\hbox{\char`\,}}{\Wrappedafterbreak}{\hbox{\char`\,}}}% 
            \lccode`\~`\;\lowercase{\def~}{\discretionary{\hbox{\char`\;}}{\Wrappedafterbreak}{\hbox{\char`\;}}}% 
            \lccode`\~`\:\lowercase{\def~}{\discretionary{\hbox{\char`\:}}{\Wrappedafterbreak}{\hbox{\char`\:}}}% 
            \lccode`\~`\?\lowercase{\def~}{\discretionary{\hbox{\char`\?}}{\Wrappedafterbreak}{\hbox{\char`\?}}}% 
            \lccode`\~`\!\lowercase{\def~}{\discretionary{\hbox{\char`\!}}{\Wrappedafterbreak}{\hbox{\char`\!}}}% 
            \lccode`\~`\/\lowercase{\def~}{\discretionary{\hbox{\char`\/}}{\Wrappedafterbreak}{\hbox{\char`\/}}}% 
            \catcode`\.\active
            \catcode`\,\active 
            \catcode`\;\active
            \catcode`\:\active
            \catcode`\?\active
            \catcode`\!\active
            \catcode`\/\active 
            \lccode`\~`\~ 	
        }
    \makeatother

    \let\OriginalVerbatim=\Verbatim
    \makeatletter
    \renewcommand{\Verbatim}[1][1]{%
        %\parskip\z@skip
        \sbox\Wrappedcontinuationbox {\Wrappedcontinuationsymbol}%
        \sbox\Wrappedvisiblespacebox {\FV@SetupFont\Wrappedvisiblespace}%
        \def\FancyVerbFormatLine ##1{\hsize\linewidth
            \vtop{\raggedright\hyphenpenalty\z@\exhyphenpenalty\z@
                \doublehyphendemerits\z@\finalhyphendemerits\z@
                \strut ##1\strut}%
        }%
        % If the linebreak is at a space, the latter will be displayed as visible
        % space at end of first line, and a continuation symbol starts next line.
        % Stretch/shrink are however usually zero for typewriter font.
        \def\FV@Space {%
            \nobreak\hskip\z@ plus\fontdimen3\font minus\fontdimen4\font
            \discretionary{\copy\Wrappedvisiblespacebox}{\Wrappedafterbreak}
            {\kern\fontdimen2\font}%
        }%
        
        % Allow breaks at special characters using \PYG... macros.
        \Wrappedbreaksatspecials
        % Breaks at punctuation characters . , ; ? ! and / need catcode=\active 	
        \OriginalVerbatim[#1,codes*=\Wrappedbreaksatpunct]%
    }
    \makeatother

    % Exact colors from NB
    \definecolor{incolor}{HTML}{303F9F}
    \definecolor{outcolor}{HTML}{D84315}
    \definecolor{cellborder}{HTML}{CFCFCF}
    \definecolor{cellbackground}{HTML}{F7F7F7}
    
    % prompt
    \makeatletter
    \newcommand{\boxspacing}{\kern\kvtcb@left@rule\kern\kvtcb@boxsep}
    \makeatother
    \newcommand{\prompt}[4]{
        {\ttfamily\llap{{\color{#2}[#3]:\hspace{3pt}#4}}\vspace{-\baselineskip}}
    }
    

    
    % Prevent overflowing lines due to hard-to-break entities
    \sloppy 
    % Setup hyperref package
    \hypersetup{
      breaklinks=true,  % so long urls are correctly broken across lines
      colorlinks=true,
      urlcolor=urlcolor,
      linkcolor=linkcolor,
      citecolor=citecolor,
      }
    % Slightly bigger margins than the latex defaults
    
    \geometry{verbose,tmargin=1in,bmargin=1in,lmargin=1in,rmargin=1in}
    
    

\begin{document}
    
    \maketitle
    
    

    
    \section{Электродинамический расчёт волнового двигателя с внутренним
расходом энергии Ф.Ф.Менде
(MenDrive)}\label{ux44dux43bux435ux43aux442ux440ux43eux434ux438ux43dux430ux43cux438ux447ux435ux441ux43aux438ux439-ux440ux430ux441ux447ux451ux442-ux432ux43eux43bux43dux43eux432ux43eux433ux43e-ux434ux432ux438ux433ux430ux442ux435ux43bux44f-ux441-ux432ux43dux443ux442ux440ux435ux43dux43dux438ux43c-ux440ux430ux441ux445ux43eux434ux43eux43c-ux44dux43dux435ux440ux433ux438ux438-ux444.ux444.ux43cux435ux43dux434ux435-mendrive}

А.Ю.Дроздов

    рассчитаем тягу в

http://fmnauka.narod.ru/dvigatel\_emdrive.pdf

Ф.Ф. МЕНДЕ, ВОЛНОВОЙ ДВИГАТЕЛЬ С ВНУТРЕННИМ РАСХОДОМ ЭНЕРГИИ
ЭЛЕКТРОМАГНИТНЫХ КОЛЕБАНИЙ

    геометрия задачи следующая:

поверхность хорошего проводника \(x<=-a\)

вакуум внутри резонатора заполняет область \(-a<=x<=a\)

В области \(x>a\) - плохой проводник

    Итак начинаем расчёт с записи первой пары уравнений Максвелла

    \[\begin{array}{c|c} \\
rot\,\vec{E} = -\frac{1}{c}\frac{\partial \vec B}{\partial t} - \frac{4\pi}{c}\vec j_m&
rot\,\vec{E} = -\frac{\partial \vec B}{\partial t} - \vec j_m \\
rot\,\vec{H} = \frac{1}{c}\frac{\partial \vec D}{\partial t} + \vec j_e &
rot\,\vec{H} = \frac{\partial \vec D}{\partial t} + \vec j_e \\
\\ \hline
\hline \\
\vec{D} = \epsilon \vec{E}&
\vec{D} = \epsilon \epsilon_0 \vec{E} \\
\vec {B} = \mu \vec{H} &
\vec {B} = \mu \mu_0 \vec{H} \\
\end{array}\]

    ищем решение в виде позволяющем разделить переменные

\[\vec{E}(t) = \vec{E}(x, y)\,e^{i\,K_z\,z}\,e^{-i\,\omega\,t}\]

\[\vec{E} = \vec{E}(x, y)\,e^{i\,K_z\,z}\]

    \[\vec{H}(t) = \vec{H}(x, y)\,e^{i\,K_z\,z}\,e^{-i\,\omega\,t}\]

\[\vec{H} = \vec{H}(x, y)\,e^{i\,K_z\,z}\]

    \[\begin{array}{c|c} \\
rot\,\vec{E} = -\frac{1}{c}\frac{\partial \mu \vec H}{\partial t} - \frac{4\pi}{c}\vec j_m = \frac{i\,\omega}{c}\,\mu\,\vec{H} - \frac{4\pi}{c}\sigma_m \vec H &
rot\,\vec{E} = -\frac{\partial \mu  \mu_0 \vec H}{\partial t} - \vec j_m = i\,\omega\,\mu\,\mu_0\,\vec{H} - \sigma_m \vec H \\
rot\,\vec{H} = \frac{1}{c}\frac{\partial \epsilon \vec E}{\partial t} + \frac{4\pi}{c} \vec j_e = - \frac{i\,\omega}{c} \epsilon \vec{E} + \frac{4\pi}{c}\sigma_e \vec E &
rot\,\vec{H} = \frac{\partial \epsilon \epsilon_0 \vec E}{\partial t} + \vec j_e = - i\,\omega \epsilon \epsilon_0 \vec{E} + \sigma_e \vec E \\
\end{array}\]

    \subsection{Аналитический вывод формул пондеромоторных
сил}\label{ux430ux43dux430ux43bux438ux442ux438ux447ux435ux441ux43aux438ux439-ux432ux44bux432ux43eux434-ux444ux43eux440ux43cux443ux43b-ux43fux43eux43dux434ux435ux440ux43eux43cux43eux442ux43eux440ux43dux44bux445-ux441ux438ux43b}

    посчитаем:

\begin{enumerate}
\def\labelenumi{\arabic{enumi})}
\item
  пондеромоторную силу, приложенную к хорошему проводнику
\item
  пондеромоторную силу, приложенную к магнетику типа феррита
\end{enumerate}

    И.Е.Тамм, Основы теории электричества, М. 1957, стр. 94

Пример. Определить силы, действующие на твердый диполь исходя из его
энергии во внешнем поле. Энергия диполя во внешнем поле, согласно
(15.8), равна

\[W = -\left(\vec {p} \vec{E} \right) = -p E cos \theta\]

Равнодействующая, приложенная к диполю сил

\[F = - \nabla W = \nabla \left(\vec {p} \vec{E} \right)\]

\[F_x = -\frac{\partial W}{\partial x}\]

для всякого не зависящего от координат вектора \(\vec{p}\)

\[\nabla \left(\vec {p} \vec{E} \right) = \left(\vec {p} \nabla \right) \cdot \vec {E} + \left[\vec{p} \times rot \vec{E}\right]\]

чтобы доказать это соотношение достаточно рассмотреть слагающую вектора
\(\nabla \left(\vec {p} \vec{E} \right) - \left(\vec {p} \nabla \right) \cdot \vec {E}\)
хотя быпо оси \(x\).

\[\begin{array}{c} \\
\frac{\partial}{\partial x} \left(\vec {p} \vec{E} \right) - \left(\vec {p} \nabla \right) \cdot E_x = \\
= p_x \frac{\partial E_x}{\partial x} + p_y \frac{\partial E_y}{\partial x} + p_z \frac{\partial E_z}{\partial x} - \left(p_x \frac{\partial E_x}{\partial x} + p_y \frac{\partial E_x}{\partial y} + p_z \frac{\partial E_x}{\partial z}\right) = \\
= p_y \left( \frac{\partial E_y}{\partial x} - \frac{\partial E_x}{\partial y} \right) 
+ p_z \left( \frac{\partial E_z}{\partial x} - \frac{\partial E_x}{\partial z} \right) =\\
= p_y rot_z E - p_z rot_y E = \left[\vec {p} \times rot \vec {E}\right]_x
\end{array}\]

    Плотность сил, приложенных к диэлектрику

\[
f = \left(\vec {P} \nabla \right) \cdot \vec {E} + \left[\vec{P} \times rot \vec{E}\right]
\]

    \[
f_x
= P_x \frac{\partial E_x}{\partial x}
+ P_y \frac{\partial E_y}{\partial x}
+ P_z \frac{\partial E_z}{\partial x}
+\left[\vec{P} \times rot \vec{E}\right]_x
\]

    где поляризация диэлектрика равна

\[ \vec{P} = \frac{\epsilon - 1}{4 \pi}\vec{E} \]

поэтому

\[
f = \left(\frac{\epsilon - 1}{4 \pi}\vec{E} \nabla \right) \cdot \vec {E} + \left[\frac{\epsilon - 1}{4 \pi}\vec{E} \times rot \vec{E}\right]
\]

    \[
f_x
= \left(\frac{\epsilon - 1}{4 \pi}\vec{E}\right)_x \frac{\partial E_x}{\partial x} \\
+ \left(\frac{\epsilon - 1}{4 \pi}\vec{E}\right)_y \frac{\partial E_y}{\partial x} \\
+ \left(\frac{\epsilon - 1}{4 \pi}\vec{E}\right)_z \frac{\partial E_z}{\partial x} \\
+ \left[\left(\frac{\epsilon - 1}{4 \pi}\vec{E}\right) \times rot \vec{E}\right]_x
\]

    В случае диагональности тензора диэлектрической проницаемости

\[
f_x
= \left(\frac{\epsilon_{xx} - 1}{4 \pi}E_x\right) \frac{\partial E_x}{\partial x}
+ \left(\frac{\epsilon_{yy} - 1}{4 \pi}E_y\right) \frac{\partial E_y}{\partial x}
+ \left(\frac{\epsilon_{zz} - 1}{4 \pi}E_z\right) \frac{\partial E_z}{\partial x}
+ \left[\left(\frac{\epsilon - 1}{4 \pi}\vec{E}\right) \times rot \vec{E}\right]_x
\]

при интегрировании по глубине материала \(x\)

    \[
\int f_x dx
= \left(\frac{\epsilon_{xx} - 1}{4 \pi}E_x\right) \frac{ E_x}{2} \Bigg|_{a}^{\infty} \\
+ \left(\frac{\epsilon_{yy} - 1}{4 \pi}E_y\right) \frac{ E_y}{2} \Bigg|_{a}^{\infty} \\
+ \left(\frac{\epsilon_{zz} - 1}{4 \pi}E_z\right) \frac{ E_z}{2} \Bigg|_{a}^{\infty} \\
+ \int \left[\left(\frac{\epsilon - 1}{4 \pi}\vec{E}\right) \times rot \vec{E}\right]_x dx
\]

при условии полного угасания поля в скин слое первые три слагаемые
приводят к формулам давления создаваемым электрическим полем из курса
электростатики Смайта

    \[\begin{array}{c} \\
\int f_x dx =
- \left(\frac{\epsilon_{xx} - 1}{4 \pi}E_x(x=a)\right) \frac{ E_x(x=a)}{2}
- \left(\frac{\epsilon_{yy} - 1}{4 \pi}E_y(x=a)\right) \frac{ E_y(x=a)}{2} \\
- \left(\frac{\epsilon_{zz} - 1}{4 \pi}E_z(x=a)\right) \frac{ E_z(x=a)}{2}
+ \int \left[\left(\frac{\epsilon - 1}{4 \pi}\vec{E}\right) \times rot \vec{E}\right]_x dx
\end{array}\]

    \[\begin{array}{c} \\
\int f_x dx
= \left(\frac{1 - \epsilon_{xx}}{4 \pi}E_x(x=a)\right) \frac{ E_x(x=a)}{2}
+ \left(\frac{1 - \epsilon_{yy}}{4 \pi}E_y(x=a)\right) \frac{ E_y(x=a)}{2} \\
+ \left(\frac{1 - \epsilon_{zz}}{4 \pi}E_z(x=a)\right) \frac{ E_z(x=a)}{2}
+ \int \left[\left(\frac{\epsilon - 1}{4 \pi}\vec{E}\right) \times rot \vec{E}\right]_x dx
\end{array}\]

    \[rot\,\vec{E} = -\frac{1}{c}\frac{\partial \vec B}{\partial t} - \frac{4\pi}{c}\vec j_m = \mu\,\frac{i\,\omega}{c}\,\vec{H} - \frac{4\pi}{c}\sigma_m \vec H \]

    \[rot\,\vec{E} = -\frac{1}{c}\frac{\partial \vec B}{\partial t} - \frac{4\pi}{c}\vec j_m = \frac{i\,\omega}{c}\mu'\,\vec{H}\]

\[rot\,\vec{H} = \frac{1}{c}\frac{\partial \vec D}{\partial t} + \frac{4\pi}{c} \vec j_e = - \frac{i\,\omega}{c} \epsilon' \vec{E}\]

    \[\begin{array}{c} \\
f_x
= \left(\frac{\epsilon - 1}{4 \pi}\vec{E}\right)_x \frac{\partial E_x}{\partial x}
+ \left(\frac{\epsilon - 1}{4 \pi}\vec{E}\right)_y \frac{\partial E_y}{\partial x} \\
+ \left(\frac{\epsilon - 1}{4 \pi}\vec{E}\right)_z \frac{\partial E_z}{\partial x}
+ \left[\left(\frac{\epsilon - 1}{4 \pi}\vec{E}\right) \times \left(\frac{i\,\omega}{c}\mu'\,\vec{H}\right)\right]_x
\end{array}\]

    И.Е.Тамм, Основы теории электричества, М. 1957, стр. 259

Потенциальная энергия магнитного диполя во внешнем поле \(H\), по
аналогии с (15.8), равна

\[U = -\left(\vec {M} \vec{H} \right) = -M H cos \theta\,\,\,\,\,(56.5)\]

Равнодействующая, приложенная к диполю, сил

\[F = - \nabla U = \nabla \left(\vec {M} \vec{H} \right)\,\,\,\,\,(56.6)\]

\[F_x = -\frac{\partial U}{\partial x}\]

И.Е.Тамм, стр. 301

Сила, действующая на систему замкнутых токов, характеризуемую магнитным
моментом \(\vec M\)

\[F = \nabla \left(\vec {M} \vec{H} \right) = \left(\vec {M} \nabla \right) \cdot \vec {H} + \left[\vec{M} \times rot \vec{H}\right] \,\,\,\,\,(66.1)\]

здесь \(\vec M\) - магнитный момент молекулы

Средняя величина силы, испытываемая отдельной молекулой магнетика

\[F = \nabla \left(\vec {M} \vec{B} \right) = \left(\vec {M} \nabla \right) \cdot \vec {B} + \left[\vec{M} \times rot \vec{B}\right] \,\,\,\,\,(66.2)\]

Плотность пондеромоторных сил, испытываемых магнетиком

\[f = \nabla \left(\vec {I} \vec{B} \right) = \left(\vec {I} \nabla \right) \cdot \vec {B} + \left[\vec{I} \times rot \vec{B}\right] \,\,\,\,\,(66.3)\]

здесь \(\vec I\) - вектор наманичения, равный магнитному моменту
молекулярных токов, рассчитанному на единицу рбъёма магнетика

\[\vec I = \sum \vec M \,\,\,\,\,(60.3)\]

\[\vec I = \kappa \vec H \,\,\,\,\,(63.1)\]

Здесь \(\kappa\) - объёмная магнитная восприимчивость

\[\mu = 1+4\pi\kappa \,\,\,\,\,(63.2)\]

\[\kappa  = \frac{\mu - 1}{4\pi}\]

если по рассматриваемому объёму протекут помимо токов молекулярных ещё и
токи проводимости, то

\[f = \frac{1}{c} \left[\vec {j_e} \times \vec{B} \right] + \left(\vec {I} \nabla \right) \cdot \vec {B} + \left[\vec{I} \times rot \vec{B}\right] \,\,\,\,\,(66.4)\]

Выражая \(\vec I = \frac{\mu - 1}{4\pi} \vec H\)

\[f = \frac{1}{c} \left[\vec {j_e} \times \vec{B} \right] + \left(\frac{\mu - 1}{4\pi} \vec H \nabla \right) \cdot \vec {B} + \left[\left(\frac{\mu - 1}{4\pi} \vec H \right) \times rot \vec{B}\right] \,\,\,\,\,(66.4)\]

    \[
f_x
= \frac{1}{c} \left[\vec {j_e} \times \vec{B} \right]_x
+ \left(\frac{\mu - 1}{4 \pi}\vec{H}\right)_x \frac{\partial B_x}{\partial x}
+ \left(\frac{\mu - 1}{4 \pi}\vec{H}\right)_y \frac{\partial B_y}{\partial x}
+ \left(\frac{\mu - 1}{4 \pi}\vec{H}\right)_z \frac{\partial B_z}{\partial x}
+ \left[\left(\frac{\mu - 1}{4 \pi}\vec{H}\right) \times rot \vec{B}\right]_x
\]

    \[\begin{array}{c} \\
f_x
= \frac{1}{c} \left[\vec {j_e} \times \vec{B} \right]_x
+ \left(\frac{\mu - 1}{4 \pi}\vec{H}\right)_x \frac{\partial \left(\mu \vec{H}\right)_x}{\partial x}
+ \left(\frac{\mu - 1}{4 \pi}\vec{H}\right)_y \frac{\partial \left(\mu \vec{H}\right)_y}{\partial x} \\
+ \left(\frac{\mu - 1}{4 \pi}\vec{H}\right)_z \frac{\partial \left(\mu \vec{H}\right)_z}{\partial x}
+ \left[\left(\frac{\mu - 1}{4 \pi}\vec{H}\right) \times rot \vec{B}\right]_x
\end{array}\]

    \[
f_x
= \left(\frac{\epsilon - 1}{4 \pi}\vec{E}\right)_x \frac{\partial E_x}{\partial x}
+ \left(\frac{\epsilon - 1}{4 \pi}\vec{E}\right)_y \frac{\partial E_y}{\partial x}
+ \left(\frac{\epsilon - 1}{4 \pi}\vec{E}\right)_z \frac{\partial E_z}{\partial x}
+ \left[\left(\frac{\epsilon - 1}{4 \pi}\vec{E}\right) \times \left(\frac{i\,\omega}{c}\mu'\,\vec{H}\right)\right]_x
\]

    \[\begin{array}{c} \\
f_x
= \frac{1}{c} \left[\vec {j_e} \times \vec{B} \right]_x
+ \left(\frac{\mu - 1}{4 \pi}\vec{H}\right)_x \frac{\partial \left(\mu \vec{H}\right)_x}{\partial x}
+ \left(\frac{\mu - 1}{4 \pi}\vec{H}\right)_y \frac{\partial \left(\mu \vec{H}\right)_y}{\partial x} \\
+ \left(\frac{\mu - 1}{4 \pi}\vec{H}\right)_z \frac{\partial \left(\mu \vec{H}\right)_z}{\partial x}
+ \left[\left(\frac{\mu - 1}{4 \pi}\vec{H}\right) \times rot \vec{B}\right]_x
\end{array}\]

    В случае диагональности тензора магнитной проницаемости

\[
f_x
= \frac{1}{c} \left[\vec {j_e} \times \vec{B} \right]_x
+ \left(\frac{\mu_{xx} - 1}{4 \pi}H_x\right) \frac{\partial H_x}{\partial x}
+ \left(\frac{\mu_{yy} - 1}{4 \pi}H_y\right) \frac{\partial H_y}{\partial x}
+ \left(\frac{\mu_{zz} - 1}{4 \pi}H_z\right) \frac{\partial H_z}{\partial x}
+ \left[\left(\frac{\mu - 1}{4 \pi}\vec{H}\right) \times rot \vec{B}\right]_x
\]

при интегрировании по глубине материала \(x\)

    \[\begin{array}{c} \\
\int f_x \, dx
= \int \frac{1}{c} \left[\vec {j_e} \times \vec{B} \right]_x \, dx
+ \left(\frac{\mu_{xx} - 1}{4 \pi}H_x\right) \frac{H_x}{2} \Bigg|_{0}^{\infty}
+ \left(\frac{\mu_{yy} - 1}{4 \pi}H_y\right) \frac{H_y}{2} \Bigg|_{0}^{\infty} \\
+ \left(\frac{\mu_{zz} - 1}{4 \pi}H_z\right) \frac{H_z}{2} \Bigg|_{0}^{\infty}
+ \int \left[\left(\frac{\mu - 1}{4 \pi}\vec{H}\right) \times rot \vec{B}\right]_x dx
\end{array}\]

при условии полного угасания поля в скин слое первые три слагаемые
приводят к магнитостатическим формулам аналогичным формулам Смайта

    Для диагонального тензора магнитной проницаемости ориентации
намагниченности \(xx\)

выражение ротора индукции магнитного поля

\[rot \vec B = \mu_{perp} rot \vec H + \left(0,\left(\frac{\mu_{par}}{\mu_{perp}}-1\right)\frac{\partial H_x}{\partial z}, 0\right)\]

\[rot \vec B = \mu_{perp} rot \vec H + \left(0,\left(\frac{\mu_{par}}{\mu_{perp}}-1\right)i\,K_z {H_x}, 0\right)\]

\[rot\,\vec{H} = - \frac{i\,\omega}{c} \, \epsilon' \vec{E}\]

\[rot \vec B = - \mu_{perp} \, \frac{i\,\omega}{c} \, \epsilon' \vec{E} + \left(0,\left(\frac{\mu_{par}}{\mu_{perp}}-1\right)i\,K_z {H_x}, 0\right)\]

%     \[\begin{array}{c} \\
% \int f_x \, dx
% = \int \frac{1}{c} \left[\vec {j_e} \times \vec{B} \right]_x \, dx
% + \left(\frac{\mu_{xx} - 1}{4 \pi}H_x\right) \frac{H_x}{2} \Bigg|_{0}^{\infty}
% + \left(\frac{\mu_{yy} - 1}{4 \pi}H_y\right) \frac{H_y}{2} \Bigg|_{0}^{\infty} \\
% + \left(\frac{\mu_{zz} - 1}{4 \pi}H_z\right) \frac{H_z}{2} \Bigg|_{0}^{\infty}
% + \int \left[\left(\frac{\mu - 1}{4 \pi}\vec{H}\right) \times rot \vec{B}\right]_x dx
% \end{array}\]

% при условии полного угасания поля в скин слое первые три слагаемые
% приводят к магнитостатическим формулам аналогичным формулам Смайта

%     \[
% \int f_x \, dx
% = \int \frac{1}{c} \left[\vec {j_e} \times \vec{B} \right]_x \, dx
% + \left(\frac{\mu_{xx} - 1}{4 \pi}H_x\right) \frac{H_x}{2} \Bigg|_{0}^{\infty}
% + \left(\frac{\mu_{yy} - 1}{4 \pi}H_y\right) \frac{H_y}{2} \Bigg|_{0}^{\infty}
% + \left(\frac{\mu_{zz} - 1}{4 \pi}H_z\right) \frac{H_z}{2} \Bigg|_{0}^{\infty}
% + \int \left[\left(\frac{\mu - 1}{4 \pi}\vec{H}\right) \times \left(-\mu_{perp} \, \frac{i\,\omega}{c} \, \epsilon' \vec{E} + \left(0,\left(\frac{\mu_{par}}{\mu_{perp}}-1\right)i\,K_z {H_x}, 0\right)\right)\right]_x dx
% \]

    Для TM волны

\[H_z = \mathit{B1}_{\mathit{zr}} e^{\left(-i \, K_{\mathit{right}_{\mathit{conductor}}} {\left(a - x\right)}\right)}\]

    \[E_y = -\mathit{B1}_{\mathit{zr}} \cdot \frac{ K_{\mathit{right}_{\mathit{conductor}}} c \mu_{r_{\mathit{xx}}} \omega e^{\left(-i \, K_{\mathit{right}_{\mathit{conductor}}} a + i \, K_{\mathit{right}_{\mathit{conductor}}} x\right)}}{c^{2} k_{z}^{2} - \epsilon_{r_{\mathit{yy}}} \mu_{r_{\mathit{xx}}} \omega^{2}}\]

    \[E_y = -\mathit{B1}_{\mathit{zr}} \cdot \frac{ K c \mu_{r_{\mathit{xx}}} \omega }{c^{2} k_{z}^{2} - \epsilon_{r_{\mathit{yy}}} \mu_{r_{\mathit{xx}}} \omega^{2}} \cdot e^{\left(-i \, K {\left(a - x\right)}\right)}\]

    Для TM волны

\[E_z = \mathit{B1}_{\mathit{zr}} e^{\left(-i \, K_{\mathit{right}_{\mathit{conductor}}} {\left(a - x\right)}\right)}\]

    \[H_y = \frac{\mathit{B1}_{\mathit{zr}} K_{\mathit{right}_{\mathit{conductor}}} c \epsilon_{r_{\mathit{xx}}} \omega e^{\left(-i \, K_{\mathit{right}_{\mathit{conductor}}} a + i \, K_{\mathit{right}_{\mathit{conductor}}} x\right)}}{c^{2} k_{z}^{2} - \epsilon_{r_{\mathit{xx}}} \mu_{r_{\mathit{yy}}} \omega^{2}}\]

    \[H_y = \frac{\mathit{B1}_{\mathit{zr}} K c \epsilon_{r_{\mathit{xx}}} \omega e^{\left(-i \, K a + i \, K x\right)}}{c^{2} k_{z}^{2} - \epsilon_{r_{\mathit{xx}}} \mu_{r_{\mathit{yy}}} \omega^{2}}\]

    \[H_y = \frac{\mathit{B1}_{\mathit{zr}} K c \epsilon_{r_{\mathit{xx}}} \omega e^{\left(-i \, K {\left(a - x\right)}\right)}}{c^{2} k_{z}^{2} - \epsilon_{r_{\mathit{xx}}} \mu_{r_{\mathit{yy}}} \omega^{2}}\]

    \[H_y = \mathit{B1}_{\mathit{zr}} \cdot \frac{ K c \epsilon_{r_{\mathit{xx}}} \omega }{c^{2} k_{z}^{2} - \epsilon_{r_{\mathit{xx}}} \mu_{r_{\mathit{yy}}} \omega^{2}} \cdot e^{\left(-i \, K {\left(a - x\right)}\right)}\]

    \[C1_{zr} = \frac{ K c \epsilon_{r_{\mathit{xx}}} \omega }{c^{2} k_{z}^{2} - \epsilon_{r_{\mathit{xx}}} \mu_{r_{\mathit{yy}}} \omega^{2}}\]

    \[E_z = \mathit{B1}_{\mathit{zr}} \cdot  e^{\left(-i \, K {\left(a - x\right)}\right)}\]

    \[H_y = \mathit{B1}_{\mathit{zr}} \cdot C1_{zr} \cdot e^{\left(-i \, K {\left(a - x\right)}\right)}\]

    Определим магнитное и электрическое поля в скин слое в упрощённом
аналитическом виде

    Здесь для сравнения не упрощённый вид этих полей полученный в ходе
решения задачи о резонаторе

    \[F_{conj} = \frac{\mathit{B1}_{\mathit{zr}} \sigma_{e_{r_{\mathit{zz}}}} \overline{\mathit{B1}_{\mathit{zr}}} \overline{\mathit{C1}_{\mathit{zr}}} e^{\left(-i \, K a + i \, a \overline{K}\right)} \int_{a}^{+\infty} e^{\left(i \, K x - i \, x \overline{K}\right)}\,{d x}}{2 \, c}\]

    \[F_{conj_l} = \frac{\mathit{B1}_{\mathit{zl}} \sigma_{e_{l_{\mathit{zz}}}} \overline{\mathit{B1}_{\mathit{zl}}} \overline{\mathit{C1}_{\mathit{zl}}} e^{\left(-i \, K_{l} a + i \, a \overline{K_{l}}\right)} \int_{-\infty}^{-a} e^{\left(-i \, K_{l} x + i \, x \overline{K_{l}}\right)}\,{d x}}{2 \, c}\]

%     \begin{tcolorbox}[breakable, size=fbox, boxrule=1pt, pad at break*=1mm,colback=cellbackground, colframe=cellborder]
% \prompt{In}{incolor}{ }{\boxspacing}
% \begin{Verbatim}[commandchars=\\\{\}]

% \end{Verbatim}
% \end{tcolorbox}

%     \[F_{conj} = \frac{\mathit{B1}_{\mathit{zr}} \sigma_{e_{r_{\mathit{zz}}}} \overline{\mathit{B1}_{\mathit{zr}}} \overline{\mathit{C1}_{\mathit{zr}}} e^{\left(-i \, K a + i \, a \overline{K}\right)} }{2 \, c} \cdot \int_{a}^{+\infty} e^{\left(i \, K x - i \, x \overline{K}\right)}\,{d x}\]

%     \[F_{conj_l} = \frac{\mathit{B1}_{\mathit{zl}} \sigma_{e_{l_{\mathit{zz}}}} \overline{\mathit{B1}_{\mathit{zl}}} \overline{\mathit{C1}_{\mathit{zl}}} e^{\left(-i \, K_{l} a + i \, a \overline{K_{l}}\right)} }{2 \, c} \cdot \int_{-\infty}^{-a} e^{\left(-i \, K_{l} x + i \, x \overline{K_{l}}\right)}\,{d x}\]

    формула давления магнитного поля на материалл в которой сдвиг фаз должен
быть учтён неявно через \(\phi_{zr} = - arg(C1_{zr})\)
\[F_{conj} = \frac{\mathit{B1}_{\mathit{zr}} \sigma_{e_{r_{\mathit{zz}}}} \overline{\mathit{B1}_{\mathit{zr}}} \overline{\mathit{C1}_{\mathit{zr}}}  }{2 \, c} \cdot e^{\left(-i \, K a + i \, a \overline{K}\right)} \cdot \int_{a}^{+\infty} e^{\left(i \, K x - i \, x \overline{K}\right)}\,{d x}\]

    \[F_{conj_l} = \frac{\mathit{B1}_{\mathit{zl}} \sigma_{e_{l_{\mathit{zz}}}} \overline{\mathit{B1}_{\mathit{zl}}} \overline{\mathit{C1}_{\mathit{zl}}}  }{2 \, c} \cdot e^{\left(-i \, K_{l} a + i \, a \overline{K_{l}}\right)} \cdot \int_{-\infty}^{-a} e^{\left(-i \, K_{l} x + i \, x \overline{K_{l}}\right)}\,{d x}\]

    Вопрос могу ли я здесь использовать формулу
\(\frac{H_{0}^{2} \cos\left(\phi\right)}{16 \, \pi} + \frac{H_{0}^{2} \sin\left(\phi\right)}{16 \, \pi}\)
пролученную в работе skin\_force\_phi если в качестве сдвига фаз
использовать \(\phi_{zr} = - arg(C1_{zr})\)

    \[F_{book} = \frac{\mathit{B1}_{\mathit{zr}}^{2} \mathit{C1}_{\mathit{zr}}^{2} }{16 \, \pi}\]

    Будет ли
\[\frac{F_{conj}}{F_{book}} = cos\left(\phi_{zr}\right) + sin\left(\phi_{zr}\right) \,\, ???\]

    Итак мы имеем определённый интеграл в выражении отношения нашей формулы
силы к книжной формуле силы

    И чтобы взять этот интеграл нужно произвести замену комплексной
переменной \(K\) на ее выражение через сумму ее действительной и мнимой
частей

    Поскольку в полученных при решении задачи о резонаторе значениях
переменных \(K\) как мнимые так и действительные части положительны

При этом нужно вспомнить что по поводу знака мнимой части переменной
\(K\) писал Зоммерфельд

\begin{verbatim}
знак перед корнем будем всегда выбирать так, чтобы корень имел положительнею мнимую часть (Зоммерфельд, Электродинамика, параграф 20, Б. Волновое поле и скин-эффект в полупространстве)
\end{verbatim}

Это как раз соответствует нашему выбору решения
$K_left_conductor_subs_d.rhs().n(), K_right_conductor_subs_d.rhs().n()$
    И поэтому мы делаем следующие assumes необходимые при дальнейшем взятии
интеграла по глубине скин слоя

    Делая подстановки и упрощая мы видим, что интеграл
\(G_{conj} = e^{\left(-i \, K a + i \, a \overline{K}\right)} \int_{a}^{+\infty} e^{\left(i \, K x - i \, x \overline{K}\right)}\,{d x}\)
взят

    С помощью такой же подстановки мы можем упростить выражение для силы
\(F_{conj} = \frac{\mathit{B1}_{\mathit{zr}} \sigma_{e_{r_{\mathit{zz}}}} \overline{\mathit{B1}_{\mathit{zr}}} \overline{\mathit{C1}_{\mathit{zr}}} e^{\left(-i \, K a + i \, a \overline{K}\right)} \int_{a}^{+\infty} e^{\left(i \, K x - i \, x \overline{K}\right)}\,{d x}}{2 \, c}\)

    Таким образом получив упрощённое аналитическое выражение для силы
давления магнитного поля на стенку нашего резонатора
\(F_{{conj}_{subs2}} = \frac{\mathit{B1}_{\mathit{zr}}^{2} \sigma_{e_{r_{\mathit{zz}}}} \overline{\mathit{C1}_{\mathit{zr}}}}{4 \, \mathit{ImK} c}\)
найдём отношение найденного только что выражения к выражению книжной
формулы \(F_{book} = \frac{H^2}{16 \pi}\)

    \begin{tcolorbox}[breakable, size=fbox, boxrule=1pt, pad at break*=1mm,colback=cellbackground, colframe=cellborder]
\prompt{In}{incolor}{252}{\boxspacing}
\begin{Verbatim}[commandchars=\\\{\}]
\PY{n}{show}\PY{p}{(}\PY{n}{F\PYZus{}book\PYZus{}r}\PY{p}{)}
\end{Verbatim}
\end{tcolorbox}

    $\displaystyle \frac{B_{H_{\mathit{zr}}} \overline{B_{H_{\mathit{zr}}}}}{16 \, \pi}$

    
    можем ли мы аналитически подобрать такие параметры резонатора для
которых влияние сдвига фаз на дисбалланс пондеромоторных сил был
максимальным

    Абрагам-подобные компоненты

    \[\begin{array}{c} \\
\int f_x \, dx
= \int \frac{1}{c} \left[\vec {j_e} \times \vec{B} \right]_x \, dx
+ \left(\frac{\mu_{xx} - 1}{4 \pi}H_x\right) \frac{H_x}{2} \Bigg|_{0}^{\infty}
+ \left(\frac{\mu_{yy} - 1}{4 \pi}H_y\right) \frac{H_y}{2} \Bigg|_{0}^{\infty} \\
+ \left(\frac{\mu_{zz} - 1}{4 \pi}H_z\right) \frac{H_z}{2} \Bigg|_{0}^{\infty}
+ \int \left[\left(\frac{\mu - 1}{4 \pi}\vec{H}\right) \times rot \vec{B}\right]_x dx
\end{array}\]

при условии полного угасания поля в скин слое первые три слагаемые
приводят к магнитостатическим формулам аналогичным формулам Смайта

Упрощая эти формулы за счёт преобразования первого интеграла, получаем
\[
\int f_x \, dx
= \int \frac{1}{c} \left[\vec {j_e} \times \vec{H} \right]_x \, dx
+ \left(\frac{\mu_{xx} - 1}{4 \pi}H_x\right) \frac{H_x}{2} \Bigg|_{0}^{\infty}
+ \int \left[\left(\frac{\mu - 1}{4 \pi}\vec{H}\right) \times rot \vec{B}\right]_x dx
\]

    моё мнение, что для устройства типа Мендрайв уравнение (105.9) Тамма не
следует распространять по всему прстранству, потому что интегральные
теоремы, типа теоремы Гаусса, как это следует из стр. 175 справочника
Корна по Математике, работают для ограниченных и пространственно
односвязных областей, причём все частные производные встречающиеся в
интегралах должны быть непрерывными. Для устройства типа Мендрайв
частная производная тензора натяжения по координате, входящая в объёмный
интеграл, терпит разрыв на границе внутреннего вакуумного промежутка со
стенками резонатора, поэтому мы вправе применить теорему Гаусса к
уравнению (105.9) не для всего пространства, а отдельно ограничившись
лишь объёмом левой стенки и отдельно ограничивший объёмом правой стенки.

Таким образом доказательство Таммом гипотезы о сохранении полного
импульса в электродинамике оказывается несостоятельным

    анализируя полученный результат на соответствие гипотезе о сохнанении
полного импульса нужно не отбрасывать как это сделал Тамм в \(\S 105\)
все дивергентные интегралы, и дело не в том что у меня в направлении
\(y\) поле не изменяется (\(k_y = 0\)), т.е. резонатор условно
бесконечен по оси \(y\), а допуская справедливость теоремы Гаусса,
преобразовавать все дивергентные интегралы в поверхностные интегралы на
внутренних поверхностях раздела металла и феррита с внутренним вакуумным
промежутком внутри стенок резонатора x=-A и x = A, учитывая, что вглубь
металла и вглубь феррита, мы считаем, что поле полностью затухает за
счёт омических или магнитных потерь

    \subsection{ГЕОМЕТРИЯ И АНАЛИЗ T\_\{xx\} ДЛЯ
TE-ВОЛНЫ}\label{ux433ux435ux43eux43cux435ux442ux440ux438ux44f-ux438-ux430ux43dux430ux43bux438ux437-t_xx-ux434ux43bux44f-te-ux432ux43eux43bux43dux44b}

\subsubsection{\texorpdfstring{\emph{Геометрия
резонатора}}{Геометрия резонатора}}\label{ux433ux435ux43eux43cux435ux442ux440ux438ux44f-ux440ux435ux437ux43eux43dux430ux442ux43eux440ux430}

\begin{equation*}
\underbrace{\text{металл}}_{x\in(-\infty,-A]}
\quad\Big|\quad
\underbrace{\text{вакуум}}_{x\in[-A,+A]}
\quad\Big|\quad
\underbrace{\text{феррит}}_{x\in[+A,+\infty)}
\end{equation*}

Поверхность \(x=-A\): граница металл/вакуум (внутренняя стенка
металла).\textbackslash{} Поверхность \(x=+A\): граница вакуум/феррит
(внутренняя стенка феррита).\textbackslash{} При \(x\to\pm\infty\): поле
затухает за счёт потерь \(\Rightarrow T_{xx}\to 0\).

    \subsubsection{\texorpdfstring{\emph{Сила на металл и феррит из потока
тензора}}{Сила на металл и феррит из потока тензора}}\label{ux441ux438ux43bux430-ux43dux430-ux43cux435ux442ux430ux43bux43b-ux438-ux444ux435ux440ux440ux438ux442-ux438ux437-ux43fux43eux442ux43eux43aux430-ux442ux435ux43dux437ux43eux440ux430}

Применяем теорему Гаусса к каждой области отдельно.

\paragraph{\texorpdfstring{Сила на металл (объём
\(x\in(-\infty,-A]\)):}{Сила на металл (объём x\textbackslash{}in(-\textbackslash{}infty,-A{]}):}}\label{ux441ux438ux43bux430-ux43dux430-ux43cux435ux442ux430ux43bux43b-ux43eux431ux44aux451ux43c-xin-infty-a}

\begin{equation*}
    F_x^{\text{мет}} = \underbrace{T_{xx}\big|_{x\to-\infty}}_{=\,0}
                     - T_{xx}\big|_{x=-A^-}
    = -T_{xx}\big|_{x=-A^-}.
\end{equation*}

\paragraph{\texorpdfstring{Сила на феррит (объём
\(x\in[+A,+\infty)\)):}{Сила на феррит (объём x\textbackslash{}in{[}+A,+\textbackslash{}infty)):}}\label{ux441ux438ux43bux430-ux43dux430-ux444ux435ux440ux440ux438ux442-ux43eux431ux44aux451ux43c-xinainfty}

\begin{equation*}
    F_x^{\text{фер}} = T_{xx}\big|_{x=+A^+}
                     - \underbrace{T_{xx}\big|_{x\to+\infty}}_{=\,0}
    = +T_{xx}\big|_{x=+A^+}.
\end{equation*}

\subsubsection{Полная тяга на
вещество:}\label{ux43fux43eux43bux43dux430ux44f-ux442ux44fux433ux430-ux43dux430-ux432ux435ux449ux435ux441ux442ux432ux43e}

\begin{equation*}
    F_x^{\text{тяга}} = F_x^{\text{мет}} + F_x^{\text{фер}}
    = T_{xx}\big|_{x=+A^+} - T_{xx}\big|_{x=-A^-}.
\end{equation*}

\subsubsection{\texorpdfstring{\emph{\(T_{xx}\) для TE-волны
(\(E_x=E_z=H_y=0\))}}{T\_\{xx\} для TE-волны (E\_x=E\_z=H\_y=0)}}\label{t_xx-ux434ux43bux44f-te-ux432ux43eux43bux43dux44b-e_xe_zh_y0}

\begin{equation*}
    T_{xx} = \frac{1}{8\pi}\bigl\{
        \mu H_x^2 - \varepsilon E_y^2 - \mu H_z^2
    \bigr\}.
\end{equation*}

\paragraph{\texorpdfstring{\emph{Граничные условия на поверхностях
раздела}}{Граничные условия на поверхностях раздела}}\label{ux433ux440ux430ux43dux438ux447ux43dux44bux435-ux443ux441ux43bux43eux432ux438ux44f-ux43dux430-ux43fux43eux432ux435ux440ux445ux43dux43eux441ux442ux44fux445-ux440ux430ux437ux434ux435ux43bux430}

На любой границе \(x = \mathrm{const}\): - \(E_y\) -\/-\/-
тангенциальная: \textbf{непрерывна} - \(H_z\) -\/-\/- тангенциальная:
\textbf{непрерывна} - \(H_x\) -\/-\/- нормальная: скачок через \(\mu\):
\(\mu_1 H_x^{(1)} = \mu_2 H_x^{(2)}\) (из непрерывности
\(B_x = \mu H_x\))

Обозначим значения поля в вакуумном зазоре у границ:

\begin{equation*}
    H_x^v \equiv H_x\big|_{\text{вак}},\quad
    H_z^v \equiv H_z\big|_{\text{вак}},\quad
    E_y^v \equiv E_y\big|_{\text{вак}}.
\end{equation*}

Тогда на стороне металла у \(x=-A\):

\begin{equation*}
    H_x^{(m)} = \frac{H_x^v}{\mu_m},\quad
    H_z^{(m)} = H_z^v,\quad
    E_y^{(m)} = E_y^v.
\end{equation*}

На стороне феррита у \(x=+A\):

\begin{equation*}
    H_x^{(f)} = \frac{H_x^v}{\mu_f},\quad
    H_z^{(f)} = H_z^v,\quad
    E_y^{(f)} = E_y^v.
\end{equation*}

\subsubsection{\texorpdfstring{\emph{Подстановка в формулу
тяги}}{Подстановка в формулу тяги}}\label{ux43fux43eux434ux441ux442ux430ux43dux43eux432ux43aux430-ux432-ux444ux43eux440ux43cux443ux43bux443-ux442ux44fux433ux438}

\begin{equation*}
    T_{xx}\big|_{x=-A^-}
    = \frac{1}{8\pi}\left\{
        \mu_m \left(\frac{H_x^v}{\mu_m}\right)^2
        - \varepsilon_m (E_y^v)^2
        - \mu_m (H_z^v)^2
    \right\}
    = \frac{1}{8\pi}\left\{
        \frac{(H_x^v)^2}{\mu_m}
        - \varepsilon_m (E_y^v)^2
        - \mu_m (H_z^v)^2
    \right\}.
\end{equation*}

\begin{equation*}
    T_{xx}\big|_{x=+A^+}
    = \frac{1}{8\pi}\left\{
        \mu_f \left(\frac{H_x^v}{\mu_f}\right)^2
        - \varepsilon_f (E_y^v)^2
        - \mu_f (H_z^v)^2
    \right\}
    = \frac{1}{8\pi}\left\{
        \frac{(H_x^v)^2}{\mu_f}
        - \varepsilon_f (E_y^v)^2
        - \mu_f (H_z^v)^2
    \right\}.
\end{equation*}

\subsection{Итоговая
тяга:}\label{ux438ux442ux43eux433ux43eux432ux430ux44f-ux442ux44fux433ux430}

\begin{equation*}
\boxed{
    F_x^{\text{тяга}} = \frac{1}{8\pi}\left[
        \left(\frac{1}{\mu_f} - \frac{1}{\mu_m}\right)(H_x^v)^2
        + (\varepsilon_m - \varepsilon_f)(E_y^v)^2
        + (\mu_m - \mu_f)(H_z^v)^2
    \right]
}
\end{equation*}

При \(\mu_f \gg 1\), \(\mu_m \approx 1\),
\(\varepsilon_m \gg \varepsilon_f\):

\begin{equation*}
    \frac{1}{\mu_f} - \frac{1}{\mu_m} \approx -1 \quad (\text{отриц., т.е. }
    \text{сила направлена в сторону металла}),
\end{equation*}\begin{equation*}
    \mu_m - \mu_f \approx -\mu_f \quad (\text{большой отрицательный вклад от }H_z).
\end{equation*}

    \begin{tcolorbox}[breakable, size=fbox, boxrule=1pt, pad at break*=1mm,colback=cellbackground, colframe=cellborder]
\prompt{In}{incolor}{ }{\boxspacing}
\begin{Verbatim}[commandchars=\\\{\}]

\end{Verbatim}
\end{tcolorbox}

    У меня в резонаторе моделируется ТE волна. Это означает ненулевой вектор
Н вдоль оси волновода, то есть вдоль оси зет, перпендикулярной рисунку.
Чтобы возбуждать такую волну можно внутрь вакуумного промежутка
поместить прямоугольный виток переменного тока.

У меня слева металл. Справа феррит.

Результат моделирования показывает, что сила действующая на металл
направлена влево, а сила действующая на феррит направлена тоже влево.

Однако природа этих сил согласно моим расчётам разная. На металл
действуют компоненты пондеромоторных сил под названием

p\_Bt\_jeB - это интеграл силы Ампера/Лоренца действующей на
тангенциальный электрический je ток Фуко возбуждаемый в скин слое
металла тангенциальным полем В, проникающим в проводник

И почему-то равная ей компонента под названием

F\_P - это интеграл сил связанных с действием градиента тангенциальной
компоненты электрического поля проникшей в скин слой проводника на
вектор электрической поляризации, который получается ненулевым за счет
того что диэлектрическая проницаемость epsilon становится комплексной за
счет слагаемого связанного с проводимостью и частотой.

Однако основная компонента сил действующих на феррит носит название

t\_I\_conv\_Hn - это натяжение за счет интеграла действия нормальной
компоненты поля H, проникшей в феррит на нормальную компоненту
намагниченности I феррита.

Если я с помощью прямоугольного витка с переменным током начну в моем
прямоугольном вакуумном промежутке резонатора возбуждать ТМ волну, то
судя по полученным мною графикам полей в вакуумном промежутке мне нужно
понять, будет ли на этот виток возбуждения действовать сила Ампера
Лоренца в магнитном поле вакуумного промежутка.

Итак, у витка 4 части.

Левая часть, параллельная оси игрек находится в области максимальной
напряжённости поля Hz недалеко от левого проводника. Исходя из общих
соображений о том, что ток возбуждения и ток Фуко обычно отталкиваются
друг от друга, то на этот участок витка будет действовать противоположно
направленная сила то есть вправо. И это наверное самая большая
компонента противосилы, которую я в своих расчётах не учитывал.

Правая часть витка, параллельная оси игрек находится в области
минимальной напряжённости поля Hz. На неё будет действовать сила
направленная влево, но уже меньшая по амплитуде.

То есть действительно аргумент о том, что противосила может быть
приложена к виткам возбуждения имеет место быть. И действительно
насколько я помню в экспериментах по Emdrive произведённой группой
немецких товарищей указывалось на то, что именно это фактор был ранее
плохо учтён в экспериментах NASA.

Но опять же хочу обратить внимание на то что основная компонента силы,
действующей на феррит в моем расчёте имеет природу, отличную от природы
силы, действующей на металл. Поэтому тут нужен особый метод расчёта,
сможет ли противосила приложенная к витку возбуждения полностью
скомпенсировать силу приложенную к ферриту.

Основная неизвестная мне величина, от которой зависит противосила, это
амплитуда переменной силы тока в витке возбуждения.

Чтобы найти эту силу тока возбуждения нужно задать следующий вопрос: а
какова формула мощности излучения прямоугольного витка возбуждения?

Приравнять эту мощность излучения мощности теплопотерь в металле и в
феррите, которая у меня уже вычисляется, найти силу тока в витке и таким
образом оценить противосилу, приложенную к витку.

    \begin{tcolorbox}[breakable, size=fbox, boxrule=1pt, pad at break*=1mm,colback=cellbackground, colframe=cellborder]
\prompt{In}{incolor}{ }{\boxspacing}
\begin{Verbatim}[commandchars=\\\{\}]

\end{Verbatim}
\end{tcolorbox}

    размер витка должен быть меньше чем 2a потому что 2a это ширина
вакуумного промежутка. а виток должен вместиться внутри вакуумного
промежутка. И у меня есть свобода выбора размежения витка внутри
промежутка. в том числе я могу разместить его несимметрично. и я могу
виток сделать более узким чем ширина зазора

Далее L - это не длина витка вдоль зет, потому зет - это ось волновода,
и для возбуждения ТМ волны плоскость витка должна быть перпендикулярна
оси зет, поэтому L - это длина витка вдоль игрек а вот с этой самой осью
игрек у меня сейчас следующая проблема. Предоставленный мною расчёт
выполнен в приближении k\_y = 0 то есть длина волновода вдоль игрек
неограничена

поэтому если оставаться в приближении k\_y = 0 то наверное у нас будет
не виток возбуждения, а длинная линия возбуждения

при этом L\_resonator - неограничен и предложенный тобой алгоритм не
работоспособен в приближении k\_y = 0

Далее

при расчёте витка возбуждения это важно

виток возбуждения имеет координату z = 0

и следовательно z = 0 - это координата максимальной амплитуды полей

при изменении координаты z поля затухают. и по хорошему нужно было бы
при расчёте сил и при расчёте мощности потерь учесть затухание полей от
z

я до сих пор неявно предполагал, что это затухание имеет однаковый
характер для силовых компонент и для мощности тепловых потерь. Так ли
это? Нужно обратить внимание на этот момент

    \begin{tcolorbox}[breakable, size=fbox, boxrule=1pt, pad at break*=1mm,colback=cellbackground, colframe=cellborder]
\prompt{In}{incolor}{ }{\boxspacing}
\begin{Verbatim}[commandchars=\\\{\}]

\end{Verbatim}
\end{tcolorbox}

    \subsubsection{СИЛА НА ВИТОК ВОЗБУЖДЕНИЯ В
TE-ВОЛНЕ}\label{ux441ux438ux43bux430-ux43dux430-ux432ux438ux442ux43eux43a-ux432ux43eux437ux431ux443ux436ux434ux435ux43dux438ux44f-ux432-te-ux432ux43eux43bux43dux435}

\subsubsection{И ЕЁ СВЯЗЬ С ТЕПЛОВЫМИ ПОТЕРЯМИ
РЕЗОНАТОРА}\label{ux438-ux435ux451-ux441ux432ux44fux437ux44c-ux441-ux442ux435ux43fux43bux43eux432ux44bux43cux438-ux43fux43eux442ux435ux440ux44fux43cux438-ux440ux435ux437ux43eux43dux430ux442ux43eux440ux430}

\section{1. Геометрия и поля
TE-волны}\label{ux433ux435ux43eux43cux435ux442ux440ux438ux44f-ux438-ux43fux43eux43bux44f-te-ux432ux43eux43bux43dux44b}

Ненулевые компоненты TE-волны в плоском резонаторе (металл \(|\) вакуум
\(|\) феррит, ось \(x\) -\/-\/- нормаль к пластинам):

\begin{equation*}
    H_x \neq 0,\quad H_z \neq 0,\quad E_y \neq 0,
    \qquad
    E_x = E_z = H_y = 0.
\end{equation*}

Поле имеет вид бегущей волны вдоль \(z\) с затуханием:

\begin{equation*}
    \mathbf{F}(x,z) = \mathbf{F}(x)\,e^{ik_z z},
    \qquad
    k_z = k_z^{\mathrm{Re}} + i k_z^{\mathrm{Im}},\quad k_z^{\mathrm{Im}} > 0.
\end{equation*}

    \subsection{2. Виток возбуждения как длинная линия
тока}\label{ux432ux438ux442ux43eux43a-ux432ux43eux437ux431ux443ux436ux434ux435ux43dux438ux44f-ux43aux430ux43a-ux434ux43bux438ux43dux43dux430ux44f-ux43bux438ux43dux438ux44f-ux442ux43eux43aux430}

Виток -\/-\/- прямоугольный контур с проводниками вдоль \(y\):

\begin{itemize}
\tightlist
\item
  левый провод: \(x = x_1\) (ближе к металлу), ток \(+I_y\,\hat{y}\)
\item
  правый провод: \(x = x_2\) (ближе к ферриту), ток \(-I_y\,\hat{y}\)
\end{itemize}

Оба провода находятся в вакуумном зазоре: \(x_1, x_2 \in (-A, +A)\).

    \subsection{3. Сила Ампера на каждый
провод}\label{ux441ux438ux43bux430-ux430ux43cux43fux435ux440ux430-ux43dux430-ux43aux430ux436ux434ux44bux439-ux43fux440ux43eux432ux43eux434}

Сила на элемент тока \(I_y\,dy\,\hat{y}\) в поле
\(\mathbf{B} = \mu_0 \mathbf{H}\) (в вакуумном зазоре \(\mu = 1\), в
системе СГС \(c\) вместо \(\mu_0\)):

\begin{equation*}
    d\mathbf{F} = \frac{1}{c}\,I_y\,dy\;(\hat{y} \times \mathbf{B}).
\end{equation*}

Для TE-волны \(B_z = H_z\), \(B_x = H_x\), \(B_y = 0\), поэтому:

\begin{equation*}
    \hat{y} \times \mathbf{B}
    = \hat{y} \times (H_x\,\hat{x} + H_z\,\hat{z})
    = H_x\,(\hat{y}\times\hat{x}) + H_z\,(\hat{y}\times\hat{z})
    = -H_x\,\hat{z} + H_z\,\hat{x}.
\end{equation*}

Сила на единицу длины по \(y\) вдоль оси \(x\):

\subsubsection{\texorpdfstring{Левый провод (\(x = x_1\), ток
\(+I_y\)):}{Левый провод (x = x\_1, ток +I\_y):}}\label{ux43bux435ux432ux44bux439-ux43fux440ux43eux432ux43eux434-x-x_1-ux442ux43eux43a-i_y}

\begin{equation*}
    \frac{F_x^{\mathrm{left}}}{L_y}
    = +\frac{I_y}{c}\,H_z(x_1).
\end{equation*}

\subsubsection{\texorpdfstring{Правый провод (\(x = x_2\), ток
\(-I_y\)):}{Правый провод (x = x\_2, ток -I\_y):}}\label{ux43fux440ux430ux432ux44bux439-ux43fux440ux43eux432ux43eux434-x-x_2-ux442ux43eux43a--i_y}

\begin{equation*}
    \frac{F_x^{\mathrm{right}}}{L_y}
    = -\frac{I_y}{c}\,H_z(x_2).
\end{equation*}

\subsubsection{\texorpdfstring{Суммарная противосила на виток вдоль
\(x\):}{Суммарная противосила на виток вдоль x:}}\label{ux441ux443ux43cux43cux430ux440ux43dux430ux44f-ux43fux440ux43eux442ux438ux432ux43eux441ux438ux43bux430-ux43dux430-ux432ux438ux442ux43eux43a-ux432ux434ux43eux43bux44c-x}

\begin{equation}
    \boxed{
        \frac{F_x^{\mathrm{wire}}}{L_y}
        = \frac{I_y}{c}\bigl[H_z(x_1) - H_z(x_2)\bigr].
    }
    \label{eq:F_wire}
\end{equation}

Знак: если \(H_z(x_1) > H_z(x_2)\) (поле больше у металла, что типично
для ветви 5), то \(F_x^{\mathrm{wire}} > 0\) -\/-\/- направлена вправо,
то есть \emph{против тяги} (тяга направлена влево, к металлу).

\subsubsection{\texorpdfstring{Замечание о силе вдоль
\(z\):}{Замечание о силе вдоль z:}}\label{ux437ux430ux43cux435ux447ux430ux43dux438ux435-ux43e-ux441ux438ux43bux435-ux432ux434ux43eux43bux44c-z}

\begin{equation*}
    \frac{F_z^{\mathrm{wire}}}{L_y}
    = -\frac{I_y}{c}\bigl[H_x(x_1) - H_x(x_2)\bigr].
\end{equation*}

Эта компонента направлена вдоль оси волновода и не связана с тягой
поперёк пластин.

    \begin{tcolorbox}[breakable, size=fbox, boxrule=1pt, pad at break*=1mm,colback=cellbackground, colframe=cellborder]
\prompt{In}{incolor}{ }{\boxspacing}
\begin{Verbatim}[commandchars=\\\{\}]

\end{Verbatim}
\end{tcolorbox}

    \subsection{4. Мощность возбуждения и связь с тепловыми
потерями}\label{ux43cux43eux449ux43dux43eux441ux442ux44c-ux432ux43eux437ux431ux443ux436ux434ux435ux43dux438ux44f-ux438-ux441ux432ux44fux437ux44c-ux441-ux442ux435ux43fux43bux43eux432ux44bux43cux438-ux43fux43eux442ux435ux440ux44fux43cux438}

\subsubsection{4.1. Мощность излучения одного
провода}\label{ux43cux43eux449ux43dux43eux441ux442ux44c-ux438ux437ux43bux443ux447ux435ux43dux438ux44f-ux43eux434ux43dux43eux433ux43e-ux43fux440ux43eux432ux43eux434ux430}

Рассмотрим один бесконечный прямолинейный провод вдоль \(y\) в точке
\(x = x_0\) с линейной плотностью тока \(I_y\) (А/см в СГС).

Работа электрического поля на токе за единицу времени на единицу длины
по \(y\) -\/-\/- это мощность, отдаваемая проводом полю:

\begin{equation*}
    \frac{P(x_0)}{L_y}
    = \frac{1}{2}\,\mathrm{Re}\bigl[E_y(x_0)\cdot I_y^*\bigr].
\end{equation*}

\emph{Обоснование:} сила, действующая на заряд \(dq\) в поле \(E_y\),
совершает работу \(dA = E_y\,dq\). Для тока \(I_y = dq/dt\) на длине
\(dy\) мощность \(dP = E_y\,I_y\,dy\). После усреднения по времени для
комплексных амплитуд получаем стандартное выражение через
\(\frac{1}{2}\mathrm{Re}[E_y I_y^*]\).

Здесь \(E_y(x_0)\) -\/-\/- это \emph{полное} поле в точке \(x_0\),
включающее как поле самого провода, так и поле резонатора. Однако поле
самого провода -\/-\/- чисто реактивное (ближнее поле бесконечного
провода), и его вклад в \(\mathrm{Re}[E_y I_y^*]\) равен нулю. Поэтому в
формуле мощности участвует только поле резонаторной моды:

\begin{equation*}
    \frac{P(x_0)}{L_y}
    = \frac{1}{2}\,\mathrm{Re}\bigl[E_y^{\mathrm{mode}}(x_0)\cdot I_y^*\bigr].
\end{equation*}

\subsubsection{4.2. Два провода витка излучают разную
мощность}\label{ux434ux432ux430-ux43fux440ux43eux432ux43eux434ux430-ux432ux438ux442ux43aux430-ux438ux437ux43bux443ux447ux430ux44eux442-ux440ux430ux437ux43dux443ux44e-ux43cux43eux449ux43dux43eux441ux442ux44c}

Виток состоит из двух проводов:

\begin{itemize}
\tightlist
\item
  левый провод в точке \(x_1\), ток \(+I_y\)
\item
  правый провод в точке \(x_2\), ток \(-I_y\)
\end{itemize}

Поскольку \(E_y^{\mathrm{mode}}(x_1) \neq E_y^{\mathrm{mode}}(x_2)\)
(поле моды неоднородно по \(x\)), мощности двух проводов различны:

\begin{equation*}
    \frac{P_1}{L_y}
    = \frac{1}{2}\,\mathrm{Re}\bigl[E_y(x_1)\cdot I_y^*\bigr],
\end{equation*}\begin{equation*}
    \frac{P_2}{L_y}
    = \frac{1}{2}\,\mathrm{Re}\bigl[E_y(x_2)\cdot(-I_y^*)\bigr]
    = -\frac{1}{2}\,\mathrm{Re}\bigl[E_y(x_2)\cdot I_y^*\bigr].
\end{equation*}

\paragraph{Знак минус у правого провода -\/-\/- потому что
ток}\label{ux437ux43dux430ux43a-ux43cux438ux43dux443ux441-ux443-ux43fux440ux430ux432ux43eux433ux43e-ux43fux440ux43eux432ux43eux434ux430-----ux43fux43eux442ux43eux43cux443-ux447ux442ux43e-ux442ux43eux43a}

в нём направлен противоположно: \(-I_y\).

Суммарная мощность, излучаемая витком:

\begin{equation}
    \boxed{
        \frac{P_{\mathrm{total}}}{L_y}
        = \frac{P_1 + P_2}{L_y}
        = \frac{1}{2}\,\mathrm{Re}\bigl[
            \bigl(E_y(x_1) - E_y(x_2)\bigr)\cdot I_y^*
          \bigr].
    }
    \label{eq:P_total}
\end{equation}

Это физически осмысленный результат: виток возбуждает поле
пропорционально \emph{разности} \(E_y\) в двух точках, то есть
пропорционально циркуляции электрического поля вдоль контура витка. Если
\(E_y(x_1) = E_y(x_2)\) (однородное поле), то суммарная мощность равна
нулю -\/-\/- виток не взаимодействует с однородной модой, что физически
очевидно.

\subsubsection{4.3. Связь с ЭДС
индукции\}}\label{ux441ux432ux44fux437ux44c-ux441-ux44dux434ux441-ux438ux43dux434ux443ux43aux446ux438ux438}

Выражение \eqref{eq:P_total} можно переписать через ЭДС, наводимую модой
в контуре витка:

\begin{equation*}
    \mathcal{E}_{\mathrm{ind}}
    = \oint_L E_s\,ds
    = E_y(x_1)\cdot L_y + E_y(x_2)\cdot(-L_y)
    = \bigl(E_y(x_1) - E_y(x_2)\bigr)\cdot L_y,
\end{equation*}

откуда:

\begin{equation*}
    \frac{P_{\mathrm{total}}}{L_y}
    = \frac{1}{2}\,\mathrm{Re}\!\left[
        \frac{\mathcal{E}_{\mathrm{ind}}}{L_y}\cdot I_y^*
      \right].
\end{equation*}

Это стандартная формула мощности источника ЭДС -\/-\/- всё
самосогласованно.

\subsubsection{4.4. Баланс мощности в стационарном
режиме}\label{ux431ux430ux43bux430ux43dux441-ux43cux43eux449ux43dux43eux441ux442ux438-ux432-ux441ux442ux430ux446ux438ux43eux43dux430ux440ux43dux43eux43c-ux440ux435ux436ux438ux43cux435}

В стационарном режиме мощность, излучаемая витком, равна суммарным
тепловым потерям в стенках резонатора:

\begin{equation}
    \frac{1}{2}\,\mathrm{Re}\bigl[
        \bigl(E_y(x_1) - E_y(x_2)\bigr)\cdot I_y^*
    \bigr]
    = w_{\mathrm{total}},
    \label{eq:balance}
\end{equation}

где \(w_{\mathrm{total}}\) -\/-\/- суммарные потери в металле и феррите
на единицу длины по \(y\) (в нормировке \texttt{calc\_thrust()}).

Записывая \(I_y = |I_y|\,e^{i\varphi}\) и вводя разность полей:

\begin{equation*}
    \Delta E_y \equiv E_y(x_1) - E_y(x_2),
\end{equation*}

уравнение баланса принимает вид:

\begin{equation*}
    \frac{1}{2}\,|I_y|\,\mathrm{Re}\bigl[\Delta E_y\,e^{-i\varphi}\bigr]
    = w_{\mathrm{total}}.
\end{equation*}

\emph{Оптимальная фаза тока} (максимальная передача мощности при
заданном \(|I_y|\)):

\begin{equation*}
    \varphi_{\mathrm{opt}} = \arg(\Delta E_y),
\end{equation*}

при которой \(\mathrm{Re}[\Delta E_y\,e^{-i\varphi}] = |\Delta E_y|\).

\emph{Амплитуда тока возбуждения:}

\begin{equation}
    \boxed{
        |I_y| = \frac{2\,w_{\mathrm{total}}}{|\Delta E_y|}
              = \frac{2\,w_{\mathrm{total}}}{|E_y(x_1) - E_y(x_2)|}.
    }
    \label{eq:I_y_correct}
\end{equation}

\subsubsection{4.5. Физический смысл: почему провода излучают
по-разному\}}\label{ux444ux438ux437ux438ux447ux435ux441ux43aux438ux439-ux441ux43cux44bux441ux43b-ux43fux43eux447ux435ux43cux443-ux43fux440ux43eux432ux43eux434ux430-ux438ux437ux43bux443ux447ux430ux44eux442-ux43fux43e-ux440ux430ux437ux43dux43eux43cux443}

Поле моды \(E_y(x)\) в вакуумном зазоре -\/-\/- это суперпозиция
экспоненциальных функций \(e^{\pm k_x x}\), убывающих от каждой стенки.
Поэтому \(|E_y(x_1)| \neq |E_y(x_2)|\) в общем случае.

Физически: левый провод (у металла) находится в более сильном
электрическом поле, если мода сосредоточена у металлической стенки, и
излучает больше мощности. Правый провод (у феррита) -\/-\/- в более
слабом поле, излучает меньше. Их вклады частично компенсируются, и
суммарная мощность определяется разностью \(\Delta E_y\).

Это аналогично тому, как сила на виток определяется разностью
\(H_z(x_1) - H_z(x_2)\): оба эффекта (сила и мощность) определяются
\emph{неоднородностью} поля по \(x\) в вакуумном зазоре.

\subsubsection{4.6. Исправленная формула удельной
противосилы}\label{ux438ux441ux43fux440ux430ux432ux43bux435ux43dux43dux430ux44f-ux444ux43eux440ux43cux443ux43bux430-ux443ux434ux435ux43bux44cux43dux43eux439-ux43fux440ux43eux442ux438ux432ux43eux441ux438ux43bux44b}

Подставляя \eqref{eq:I_y_correct} в формулу силы на виток:

\begin{equation*}
    \frac{F_x^{\mathrm{wire}}}{L_y}
    = \frac{|I_y|}{c}\bigl[H_z(x_1) - H_z(x_2)\bigr]
    = \frac{2\,w_{\mathrm{total}}}{c}\cdot
      \frac{H_z(x_1) - H_z(x_2)}{|E_y(x_1) - E_y(x_2)|}.
\end{equation*}

Удельная противосила (на единицу рассеиваемой мощности):

\begin{equation}
    \boxed{
        \frac{F_x^{\mathrm{wire}}/L_y}{w_{\mathrm{total}}}
        = \frac{2}{c}\cdot
          \frac{H_z(x_1) - H_z(x_2)}{|E_y(x_1) - E_y(x_2)|}.
    }
    \label{eq:specific_counter_correct}
\end{equation}

Отношение \(\dfrac{H_z(x_1) - H_z(x_2)}{|E_y(x_1) - E_y(x_2)|}\) имеет
размерность \(1/Z_{\mathrm{eff}}\), где \(Z_{\mathrm{eff}}\) -\/-\/-
эффективное волновое сопротивление моды, ``видимое'' витком возбуждения.

При \(Z_{\mathrm{eff}} \gg 1\) (высокое волновое сопротивление):
противосила \(\ll\) тяга, нетто-тяга велика.

При \(Z_{\mathrm{eff}} \sim 1\) (нормальное волновое сопротивление):
противосила \(\sim\) тяга, нетто-тяга мала.

    \subsection{5. Удельная противосила на
виток}\label{ux443ux434ux435ux43bux44cux43dux430ux44f-ux43fux440ux43eux442ux438ux432ux43eux441ux438ux43bux430-ux43dux430-ux432ux438ux442ux43eux43a}

Подставляя исправленную амплитуду тока I\_y\_correct в формулу суммарной
силы на виток F\_wire:

\begin{equation*}
    \frac{F_x^{\mathrm{wire}}}{L_y}
    = \frac{|I_y|}{c}\bigl[H_z(x_1) - H_z(x_2)\bigr]
    = \frac{2\,w_{\mathrm{total}}}{c}\cdot
      \frac{H_z(x_1) - H_z(x_2)}{|E_y(x_1) - E_y(x_2)|}.
\end{equation*}

Удельная противосила (на единицу рассеиваемой мощности):

\begin{equation}
    \boxed{
        \frac{F_x^{\mathrm{wire}}/L_y}{w_{\mathrm{total}}}
        = \frac{2}{c}\cdot
          \frac{H_z(x_1) - H_z(x_2)}{|E_y(x_1) - E_y(x_2)|}
        \equiv \frac{2}{c\,Z_{\mathrm{eff}}},
    }
    \label{eq:specific_counter}
\end{equation}

где введено эффективное волновое сопротивление моды, "видимое" витком
возбуждения:

\begin{equation}
    Z_{\mathrm{eff}}
    \equiv \frac{|E_y(x_1) - E_y(x_2)|}{H_z(x_1) - H_z(x_2)}.
    \label{eq:Z_eff}
\end{equation}

\emph{Замечание о знаках:} формула записана для случая
\(H_z(x_1) > H_z(x_2)\) (поле \(H_z\) больше у металлической стенки, что
типично для ветви 5). В общем случае нужно брать \(\mathrm{Re}[\ldots]\)
аналогично формуле мощности.

    \subsection{6. Степень компенсации
тяги}\label{ux441ux442ux435ux43fux435ux43dux44c-ux43aux43eux43cux43fux435ux43dux441ux430ux446ux438ux438-ux442ux44fux433ux438}

Пусть \(\mathcal{T}\) {[}Н/Вт{]} -\/-\/- удельная тяга на вещество
(металл + феррит), вычисленная в \texttt{calc\_thrust()}. Степень
компенсации тяги противосилой витка:

\begin{equation}
    \boxed{
        \eta \equiv
        \frac{F_x^{\mathrm{wire}}/L_y}{w_{\mathrm{total}}\cdot\mathcal{T}}
        = \frac{2}{c\,Z_{\mathrm{eff}}\cdot\mathcal{T}}.
    }
    \label{eq:eta}
\end{equation}

Нетто-тяга резонатора на единицу мощности:

\begin{equation}
    \mathcal{T}_{\mathrm{net}}
    = \mathcal{T} - \frac{F_x^{\mathrm{wire}}/L_y}{w_{\mathrm{total}}}
    = \mathcal{T}\,(1 - \eta)
    = \mathcal{T}\!\left(1 - \frac{2}{c\,Z_{\mathrm{eff}}\cdot\mathcal{T}}\right).
    \label{eq:T_net}
\end{equation}

Предельные случаи: - \(Z_{\mathrm{eff}} \gg \dfrac{2}{c\,\mathcal{T}}\):
\(\eta \ll 1\), нетто-тяга \(\approx \mathcal{T}\) виток слабо
компенсирует тягу. - \(Z_{\mathrm{eff}} = \dfrac{2}{c\,\mathcal{T}}\):
\(\eta = 1\), нетто-тяга \(= 0\) полная компенсация. -
\(Z_{\mathrm{eff}} \ll \dfrac{2}{c\,\mathcal{T}}\): \(\eta \gg 1\) виток
"перекомпенсирует" тягу, что физически означает неправильный выбор
положения \(x_1\), \(x_2\) или фазы тока.

    \begin{tcolorbox}[breakable, size=fbox, boxrule=1pt, pad at break*=1mm,colback=cellbackground, colframe=cellborder]
\prompt{In}{incolor}{ }{\boxspacing}
\begin{Verbatim}[commandchars=\\\{\}]

\end{Verbatim}
\end{tcolorbox}

    \subsection{7. Физический смысл
результата}\label{ux444ux438ux437ux438ux447ux435ux441ux43aux438ux439-ux441ux43cux44bux441ux43b-ux440ux435ux437ux443ux43bux44cux442ux430ux442ux430}

\subsubsection{7.1. Импульсный баланс
системы}\label{ux438ux43cux43fux443ux43bux44cux441ux43dux44bux439-ux431ux430ux43bux430ux43dux441-ux441ux438ux441ux442ux435ux43cux44b}

Полный закон сохранения импульса для системы (металл + вакуум + феррит +
виток) при \(k_y \neq 0\):

\begin{equation*}
    F_x^{\mathrm{мет}} + F_x^{\mathrm{фер}}
    + F_x^{\mathrm{вак}}
    + F_x^{\mathrm{wire}} = 0.
\end{equation*}

При \(k_y = 0\) слагаемое \(F_x^{\mathrm{вак}}\) включает поток
электромагнитного импульса через боковые стенки \(y = \pm\infty\), и
баланс в расчёте на единицу длины по \(y\) нарушается:

\begin{equation*}
    F_x^{\mathrm{мет}} + F_x^{\mathrm{фер}} + F_x^{\mathrm{wire}}
    = -F_x^{\mathrm{вак}}
    \neq 0.
\end{equation*}

\subsubsection{7.2. Почему большая тяга совместима с малой
противосилой}\label{ux43fux43eux447ux435ux43cux443-ux431ux43eux43bux44cux448ux430ux44f-ux442ux44fux433ux430-ux441ux43eux432ux43cux435ux441ux442ux438ux43cux430-ux441-ux43cux430ux43bux43eux439-ux43fux440ux43eux442ux438ux432ux43eux441ux438ux43bux43eux439}

Из формулы \eqref{eq:specific_counter} противосила определяется
отношением \(\Delta H_z / \Delta E_y\). Большая тяга \(\mathcal{T}\)
возникает в режимах с аномально большим \(H_x\) у феррита (концентрация
магнитного потока). В этих же режимах:

\begin{itemize}
\tightlist
\item
  \(H_z\) в вакуумном зазоре может быть \emph{мало} (энергия
  сосредоточена в \(H_x\), а не в \(H_z\)), поэтому
  \(\Delta H_z = H_z(x_1) - H_z(x_2)\) мало;
\item
  \(E_y\) в вакуумном зазоре при этом конечно (из уравнения Максвелла
  \(\partial H_z/\partial x -  \partial H_x/\partial z = \frac{i\omega\varepsilon}{c}E_y\)),
  поэтому \(\Delta E_y\) не мало.
\end{itemize}

Следовательно, \(Z_{\mathrm{eff}} \gg 1\), \(\eta \ll 1\), и нетто-тяга
\(\approx \mathcal{T}\) -\/-\/- большая тяга \emph{не компенсируется}
противосилой витка.

\subsubsection{7.3. Аналогия с
трансформатором}\label{ux430ux43dux430ux43bux43eux433ux438ux44f-ux441-ux442ux440ux430ux43dux441ux444ux43eux440ux43cux430ux442ux43eux440ux43eux43c}

Виток возбуждения в резонаторе аналогичен первичной обмотке
трансформатора. Мощность передаётся через \(E_y\) (аналог напряжения), а
сила -\/-\/- через \(H_z\) (аналог тока). Если трансформатор работает в
режиме высокого напряжения / малого тока (высокое \(Z_{\mathrm{eff}}\)),
то при той же передаваемой мощности сила на обмотку мала. Именно в этом
режиме и работает ветвь 5.

    \subsection{8. Алгоритм численной
проверки}\label{ux430ux43bux433ux43eux440ux438ux442ux43c-ux447ux438ux441ux43bux435ux43dux43dux43eux439-ux43fux440ux43eux432ux435ux440ux43aux438}

\subsubsection{АНАЛИЗ ЕДИНИЦ ИЗМЕРЕНИЯ И АЛГОРИТМ ЧИСЛЕННОЙ
ПРОВЕРКИ}\label{ux430ux43dux430ux43bux438ux437-ux435ux434ux438ux43dux438ux446-ux438ux437ux43cux435ux440ux435ux43dux438ux44f-ux438-ux430ux43bux433ux43eux440ux438ux442ux43c-ux447ux438ux441ux43bux435ux43dux43dux43eux439-ux43fux440ux43eux432ux435ux440ux43aux438}

\subsection{\texorpdfstring{\emph{Анализ единиц
измерения}}{Анализ единиц измерения}}\label{ux430ux43dux430ux43bux438ux437-ux435ux434ux438ux43dux438ux446-ux438ux437ux43cux435ux440ux435ux43dux438ux44f}

\subsubsection{\texorpdfstring{Что вычисляет функция
\(\texttt{wr}/\texttt{wl}\)}{Что вычисляет функция \textbackslash{}texttt\{wr\}/\textbackslash{}texttt\{wl\}}}\label{ux447ux442ux43e-ux432ux44bux447ux438ux441ux43bux44fux435ux442-ux444ux443ux43dux43aux446ux438ux44f-textttwrtextttwl}

Функция вычисляет объёмную плотность тепловых потерь, интегрированную по
\(x\) и усреднённую по \(y\):

\begin{equation*}
    w(x) = \frac{1}{2}\Bigl(
        \sigma_e^{zz}|E_z|^2 + \sigma_e^{yy}|E_y|^2 + \sigma_e^{xx}|E_x|^2
        + \sigma_m^{zz}|H_z|^2 + \sigma_m^{yy}|H_y|^2 + \sigma_m^{xx}|H_x|^2
    \Bigr)
    \quad \left[\frac{\text{эрг}}{\text{с}\cdot\text{см}^3}\right].
\end{equation*}

После интегрирования по \(x\) (от стенки до глубины проводника) и
перевода в ватты:

\begin{equation*}
    W_{\text{total}} = \left(\int_{-A-h_l}^{-A} w_l\,dx
                      + \int_{A}^{A+h_r} w_r\,dx\right) \cdot 10^{-7}
    \quad \left[\frac{\text{Вт}}{\text{см}^2}\right].
\end{equation*}

Функция \(avg_over_y\) делит на \(L_y\), поэтому нормировка по
\(y\) уже выполнена. Интегрирование по \(z\) не делается -\/-\/- поле
берётся при \(z = 0\) с нормировкой \(e^{ik_z z}\).

\emph{Итого:} \(W_{\text{total}}\) имеет размерность
\(\left[\dfrac{\text{Вт}}{\text{см}^2}\right]\) -\/-\/- мощность потерь
на единицу площади в плоскости \((y, z)\).

    \subsection{Размерность силы и удельной
тяги}\label{ux440ux430ux437ux43cux435ux440ux43dux43eux441ux442ux44c-ux441ux438ux43bux44b-ux438-ux443ux434ux435ux43bux44cux43dux43eux439-ux442ux44fux433ux438}

Аналогично, \(\texttt{force\_N}\) -\/-\/- сила, интегрированная по \(x\)
и усреднённая по \(y\):

\begin{equation*}
    F \quad \left[\frac{\text{Н}}{\text{см}^2}\right].
\end{equation*}

Отношение:

\begin{equation*}
    \mathcal{T} = \frac{F\;[\text{Н/см}^2]}{W_{\text{total}}\;[\text{Вт/см}^2]}
    \quad \left[\frac{\text{Н}}{\text{Вт}}\right].
\end{equation*}

Единицы \(\text{см}^2\) сокращаются -\/-\/- удельная тяга в
\([\text{Н/Вт}]\) вычисляется корректно.

    \subsection{\texorpdfstring{\emph{Нормировка по \(z\) через эффективную
длину}}{Нормировка по z через эффективную длину}}\label{ux43dux43eux440ux43cux438ux440ux43eux432ux43aux430-ux43fux43e-z-ux447ux435ux440ux435ux437-ux44dux444ux444ux435ux43aux442ux438ux432ux43dux443ux44e-ux434ux43bux438ux43dux443}

Поле имеет вид \(\sim e^{ik_z z}\) с
\(k_z = k_z^{\text{Re}} + ik_z^{\text{Im}}\), поэтому интенсивность
\(\sim e^{-2k_z^{\text{Im}}|z|}\). Эффективная длина по \(z\):

\begin{equation*}
    L_z^{\text{eff}} = \int_0^{\infty} e^{-2k_z^{\text{Im}} z}\,dz
    = \frac{1}{2k_z^{\text{Im}}}
    \quad [\text{см}].
\end{equation*}

Для перехода от значения при \(z = 0\) к интегралу по \(z\) нужно
умножить на \(L_z^{\text{eff}}\). Для перехода обратно -\/-\/- разделить
на \(L_z^{\text{eff}}\).

    \subsection{\texorpdfstring{\emph{Алгоритм вычисления противосилы на
виток}}{Алгоритм вычисления противосилы на виток}}\label{ux430ux43bux433ux43eux440ux438ux442ux43c-ux432ux44bux447ux438ux441ux43bux435ux43dux438ux44f-ux43fux440ux43eux442ux438ux432ux43eux441ux438ux43bux44b-ux43dux430-ux432ux438ux442ux43eux43a}

\subsubsection{\texorpdfstring{*Шаг 1. Исходные данные из
\(\texttt{calc\_thrust()}\)}{*Шаг 1. Исходные данные из \textbackslash{}texttt\{calc\textbackslash{}\_thrust()\}}}\label{ux448ux430ux433-1.-ux438ux441ux445ux43eux434ux43dux44bux435-ux434ux430ux43dux43dux44bux435-ux438ux437-textttcalc_thrust}

\begin{align*}
    W_{\text{total}} &\quad [\text{Вт/см}^2]
        \quad\text{--- суммарные потери (металл + феррит)}, \\
    F &\quad [\text{Н/см}^2]
        \quad\text{--- суммарная тяга на вещество}, \\
    \mathcal{T} &= F / W_{\text{total}} \cdot 10^3
        \quad [\text{Н/кВт}].
\end{align*}

    \subsection{\texorpdfstring{\emph{Шаг 2. Поля в точках расположения
проводов
витка}}{Шаг 2. Поля в точках расположения проводов витка}}\label{ux448ux430ux433-2.-ux43fux43eux43bux44f-ux432-ux442ux43eux447ux43aux430ux445-ux440ux430ux441ux43fux43eux43bux43eux436ux435ux43dux438ux44f-ux43fux440ux43eux432ux43eux434ux43eux432-ux432ux438ux442ux43aux430}

Провода витка расположены в вакуумном зазоре \(x \in (-A, +A)\):

\begin{equation*}
    x_1 = -\alpha A \quad (\text{левый, ближе к металлу}),
    \qquad
    x_2 = +\alpha A \quad (\text{правый, ближе к ферриту}),
\end{equation*}

где \(\alpha \in (0, 1)\) -\/-\/- параметр положения витка.

Комплексные амплитуды полей при \(z = 0\):

\begin{equation*}
    H_z(x_1),\; H_z(x_2) \quad [\text{Гс}],
    \qquad
    E_y(x_1),\; E_y(x_2) \quad [\text{статВ/см}].
\end{equation*}

    \subsection{\texorpdfstring{\emph{Шаг 3. Разности
полей}}{Шаг 3. Разности полей}}\label{ux448ux430ux433-3.-ux440ux430ux437ux43dux43eux441ux442ux438-ux43fux43eux43bux435ux439}

\begin{equation*}
    \Delta H_z = \mathrm{Re}\bigl[H_z(x_1) - H_z(x_2)\bigr]
    \quad [\text{Гс}],
\end{equation*}

\begin{equation*}
    \Delta E_y = E_y(x_1) - E_y(x_2)
    \quad [\text{статВ/см}],
\end{equation*}

\begin{equation*}
    |\Delta E_y| = |E_y(x_1) - E_y(x_2)|.
\end{equation*}

    \subsection{\texorpdfstring{\emph{Шаг 4. Амплитуда тока
возбуждения}}{Шаг 4. Амплитуда тока возбуждения}}\label{ux448ux430ux433-4.-ux430ux43cux43fux43bux438ux442ux443ux434ux430-ux442ux43eux43aux430-ux432ux43eux437ux431ux443ux436ux434ux435ux43dux438ux44f}

Баланс мощности в согласованных единицах:

\begin{equation*}
    \underbrace{\frac{1}{2}\,\mathrm{Re}[\Delta E_y \cdot I_y^*]}_{\text{эрг/(с·см)}}
    \cdot \frac{1}{L_z^{\text{eff}}} \cdot 10^{-7}
    = W_{\text{total}} \quad [\text{Вт/см}^2].
\end{equation*}

Откуда амплитуда тока в СГС:

\begin{equation}
    \boxed{
        |I_y| = \frac{2\,W_{\text{total}}\,L_z^{\text{eff}} \cdot 10^7}{|\Delta E_y|}
        \quad [\text{статА}].
    }
    \label{eq:I_y_units}
\end{equation}

В единицах СИ (\(1\,\text{А} = 3\cdot10^9\,\text{статА}\)):

\begin{equation*}
    |I_y^{\text{SI}}| = \frac{|I_y|}{3 \cdot 10^9} \quad [\text{А}].
\end{equation*}

    \subsection{\texorpdfstring{\emph{Шаг 5. Мощности отдельных
проводов}}{Шаг 5. Мощности отдельных проводов}}\label{ux448ux430ux433-5.-ux43cux43eux449ux43dux43eux441ux442ux438-ux43eux442ux434ux435ux43bux44cux43dux44bux445-ux43fux440ux43eux432ux43eux434ux43eux432}

Мощность левого провода (ток \(+I_y\)):

\begin{equation*}
    P_1 = \frac{1}{2}\,\mathrm{Re}[E_y(x_1)\cdot I_y^*]
    \cdot \frac{10^{-7}}{L_z^{\text{eff}}}
    \quad [\text{Вт/см}^2].
\end{equation*}

Мощность правого провода (ток \(-I_y\)):

\begin{equation*}
    P_2 = -\frac{1}{2}\,\mathrm{Re}[E_y(x_2)\cdot I_y^*]
    \cdot \frac{10^{-7}}{L_z^{\text{eff}}}
    \quad [\text{Вт/см}^2].
\end{equation*}

Проверка:

\begin{equation*}
    P_1 + P_2
    = \frac{1}{2}\,\mathrm{Re}[\Delta E_y \cdot I_y^*]
    \cdot \frac{10^{-7}}{L_z^{\text{eff}}}
    = W_{\text{total}}. \quad \checkmark
\end{equation*}

    \subsubsection{\texorpdfstring{\emph{Шаг 6. Сила на
виток}}{Шаг 6. Сила на виток}}\label{ux448ux430ux433-6.-ux441ux438ux43bux430-ux43dux430-ux432ux438ux442ux43eux43a}

Сила на виток в СГС:

\begin{equation*}
    \frac{F_{\text{wire}}}{L_y \cdot L_z}
    = \frac{|I_y|}{c}\,\Delta H_z
    \quad [\text{дин/см}^2].
\end{equation*}

Нормировка на \(L_z^{\text{eff}}\) и перевод в \([\text{Н/см}^2]\)
(\(1\,\text{Н} = 10^5\,\text{дин}\)):

\begin{equation}
    \boxed{
        F_{\text{wire}}
        = \frac{|I_y|}{c}\,\frac{\Delta H_z}{L_z^{\text{eff}}} \cdot 10^{-5}
        \quad [\text{Н/см}^2].
    }
    \label{eq:F_wire_units}
\end{equation}

\subsubsection{\texorpdfstring{\emph{Шаг 7. Степень
компенсации}}{Шаг 7. Степень компенсации}}\label{ux448ux430ux433-7.-ux441ux442ux435ux43fux435ux43dux44c-ux43aux43eux43cux43fux435ux43dux441ux430ux446ux438ux438}

Удельная противосила:

\begin{equation*}
    \frac{F_{\text{wire}}}{W_{\text{total}}}
    \quad [\text{Н/Вт}].
\end{equation*}

Степень компенсации:

\begin{equation}
    \boxed{
        \eta = \frac{F_{\text{wire}} / W_{\text{total}}}{\mathcal{T} / 10^3}
             = \frac{F_{\text{wire}} \cdot 10^3}{W_{\text{total}} \cdot \mathcal{T}}.
    }
    \label{eq:eta_units}
\end{equation}

Нетто-тяга:

\begin{equation*}
    \mathcal{T}_{\text{net}} = \mathcal{T}\,(1 - \eta)
    \quad [\text{Н/кВт}].
\end{equation*}

\subsubsection{\texorpdfstring{\emph{Шаг 8. Эффективное волновое
сопротивление}}{Шаг 8. Эффективное волновое сопротивление}}\label{ux448ux430ux433-8.-ux44dux444ux444ux435ux43aux442ux438ux432ux43dux43eux435-ux432ux43eux43bux43dux43eux432ux43eux435-ux441ux43eux43fux440ux43eux442ux438ux432ux43bux435ux43dux438ux435}

\begin{equation*}
    Z_{\text{eff}} = \frac{|\Delta E_y|}{|\Delta H_z|}
    \quad \left[\frac{\text{статВ/см}}{\text{Гс}}\right]
    = \left[\frac{\text{статВ}}{\text{Гс}\cdot\text{см}}\right].
\end{equation*}

В СГС это безразмерная величина (так как \([E] = [H]\) в СГС). При
\(Z_{\text{eff}} \gg 1\): \(\eta \ll 1\), нетто-тяга велика. При
\(Z_{\text{eff}} \sim 1\): \(\eta \sim 1\), компенсация полная.

%     \begin{tcolorbox}[breakable, size=fbox, boxrule=1pt, pad at break*=1mm,colback=cellbackground, colframe=cellborder]
% \prompt{In}{incolor}{ }{\boxspacing}
% \begin{Verbatim}[commandchars=\\\{\}]

% \end{Verbatim}
% \end{tcolorbox}

%     \begin{tcolorbox}[breakable, size=fbox, boxrule=1pt, pad at break*=1mm,colback=cellbackground, colframe=cellborder]
% \prompt{In}{incolor}{ }{\boxspacing}
% \begin{Verbatim}[commandchars=\\\{\}]

% \end{Verbatim}
% \end{tcolorbox}


    % Add a bibliography block to the postdoc
    
    
    
\end{document}
